% Section 6: the vertex-on-edge law.
\section{The vertex-on-edge law}\label{sec:s7}

\begin{convention}[the standing hypotheses of this section]
\label{conv:s7standing}
Every statement of this section that speaks of a half, of a half's carriers,
rows or selectors, or of the jump those are compared against, is made under the
following three hypotheses, which are named here once and cited by name rather
than repeated or inherited silently:
\begin{enumerate}
\item[(H1)] the event is a \emph{simple vertex-on-edge event} in the sense of
Convention~\ref{conv:events}, with moving vertex $M$ and remote edge $(a,b)$,
occurring on a path inside $\mathcal R_n$ whose polygons have $n$ vertices;
\item[(H2)] the polygon $P_\bullet$ of whichever chamber is under discussion is
\emph{generic} (Definition~\ref{def:generic}(A));
\item[(H3)] both halves $\lambda_1,\lambda_2$ of Definition~\ref{def:halves},
taken at the event, are \emph{generic polygons}, each with at least three
vertices.
\end{enumerate}
When the event is of bigon type, on the side of the event where both edges
incident to $M$ cross the remote edge $(a,b)$, the two crossings so created are
the \emph{newborn} crossings: $x$, where the edge $M-1\to M$ meets $(a,b)$,
and $y$, where $M\to M+1$ meets it.
These hypotheses are in force for the whole section \emph{by this declaration},
not by each statement repeating them: a statement of this section is to be read
with (H1)--(H3) available whether or not it names them, in the way
Convention~\ref{conv:lawful} binds the letter $A$ everywhere below without
being cited at each use. Several statements do name the convention --- those
where a reader would otherwise wonder which polygons are meant --- and that is
a courtesy, not the mechanism. What the declaration is \emph{not} is an
inheritance from a neighbouring proof: the hypotheses are written out here, in
one numbered place, and a formalizer reading any statement of the section has
one address to look at.

The convention exists because the objects of this section do not exist without
it. The half quantities --- a half's carriers, its grouped rows $Q_i$, its
selectors $W_i$, its sector decomposition, and the products
$X_1(\lambda_1)X_1(\lambda_2)$ that the jump is compared against --- are read
off polygons through $X_1$, and $X_1$ is defined on \emph{generic} polygons
only (Definition~\ref{def:X1}). An earlier arrangement stated the hypothesis at
the two endpoint theorems and, in the middle of the section, said that "every
consumer below inherits" it, which is not a quantifier: a reader of the sector
definition or of the sliding restart could not see from the page what was
assumed. (H3) is discharged, at every event the induction meets, by the
certificate proved with the main theorem; nothing in this section proves it.
\end{convention}

\subsection{Grouped carriers are connected sums}

The factor $\Omega_1(S,L)$ of Definition~\ref{def:X1} reads one coefficient of
the \emph{product} $P_{S,L}=\prod_H P_H$ over the pieces carried by $L$. The
carrier floor of Theorem~\ref{thm:carrierfloor}(C), on the other hand, is a
statement about one knot diagram. The bridge is that the product is the
polynomial of a single knot.

\begin{lemma}[carrier restriction is the intrinsic piece diagram]
\label{lem:pieceintrinsic}
Let $P$ be generic, $S\in\Ind(G_P)$, and let $H$ be a residual piece of $S$,
carried by the carrier $L$ (Lemma~\ref{lem:carriers}(iv)). Let $D_L(H)$ be the
diagram obtained from $L$ by retaining exactly the crossings of $H$, each
resolved by the divide convention of Definition~\ref{def:piecediagram}, and
erasing all others; and let $D_P(H)$ be the intrinsic piece diagram of that
definition, obtained from the parent polygon $P$ by retaining exactly the same
crossings under the same convention. Throughout, that a crossing \emph{lies on}
a carrier means that both of its traversal preimages belong to that carrier,
which is the condition Lemma~\ref{lem:carriers}(iv) supplies.
Then the identity map on the visits of $H$ is a record isomorphism
(Definition~\ref{def:record}) from the record of $D_L(H)$ to the record of
$D_P(H)$. Consequently they present the same oriented link
(Axiom~\ref{ax:gausscode}), so
\[
   P_{D_L(H)}=P_H,\qquad w\bigl(D_L(H)\bigr)=|H| .
\]
\end{lemma}
\begin{proof}
\emph{Same double points.} Smoothing at a point $c\in S$ replaces the two
transversally crossing branches by two disjoint arcs inside a disc containing no
other double point, so it destroys the double point $c$ and creates none. Hence
the double points of $L$ are exactly the double points of $P$ that lie on $L$
and are not in $S$. Every label of $H$ has both of its visits on $L$ and
$H\cap S=\varnothing$, so each is a double point of $L$; and both diagrams
retain exactly those and no others.

\emph{Same over/under and signs.} The divide convention
(Definition~\ref{def:piecediagram}) reads the over/under at a double point from
the two branch directions there, and the sign of a crossing is read from the
same two directions with the orientation. Neither is changed by a smoothing
performed at a different point, which alters the curve only inside a disc
containing no other double point. Every crossing is positive in both diagrams,
so both writhes equal $|H|$.

\emph{Same cyclic order.} The traversal of $D_P(H)$ is the traversal circle
$\Gamma$ itself, so the cyclic order of the $2|H|$ visits is their order on
$\Gamma$. By Lemma~\ref{lem:carrierword} the marked points lying on $L$ are
traversed by $L$ in the order induced from $\Gamma$, and the visits of $H$ are
among them; restricting to those gives the same cyclic word.

\emph{The isomorphism.} The two records are not equal: their traversal circles
are different circles, one running around $P$ and one around $L$. What the three
paragraphs give is the map. Let $\varphi$ send each visit of $H$ on the
traversal circle of $D_L(H)$ to the same visit on the traversal circle of
$D_P(H)$ --- the same point of the plane, reached by the same branch --- which
is a bijection of marked-point sets by the first paragraph. It preserves the
cyclic order by the third, carries double points to double points because the
pairing is by label in both, and preserves the over/under designation and the
sign by the second. Those are conditions (a) to (d) of
Definition~\ref{def:record}, so $\varphi$ is a record isomorphism, and
Axiom~\ref{ax:gausscode} identifies the links.

An earlier revision of this lemma said the two records were \emph{equal}. They
are not: equality of records presupposes
one traversal circle, which is exactly what these two diagrams do not share.
Nothing in the three paragraphs changes; what changes is that the map they
produce is now named, and the definition it satisfies is printed.
\end{proof}

\begin{lemma}[HOMFLY products and the two-component $z^{-1}$ row]
\label{lem:homflyrows}
Put $R=\mathbb Z[a^{\pm1},z^{\pm1}]$ and
$\delta=(a-a^{-1})/z\in R$.
\begin{enumerate}
\item[(i)] For oriented knots $K,J$,
\[
 P_{K\# J}=P_KP_J;
\]
and for an oriented knot $K$ and an oriented link $J$,
\[
 P_{K\sqcup J}=\delta P_KP_J.
\]
\item[(ii)] Let $D$ be an oriented link diagram with ordered components
$D_1,D_2$ and linking number $\lambda$. Then
\[
 [z^{-1}]P_D=(a-a^{-1})a^{-2\lambda}
       [z^0]\bigl(P_{D_1}P_{D_2}\bigr).
\]
\item[(iii)] Consequently, for oriented knots $K_1,\dots,K_n$ ($n\geq1$),
\[
 P_{K_1\sqcup\dots\sqcup K_n}=\delta^{\,n-1}P_{K_1}\cdots P_{K_n}.
\]
\end{enumerate}
\end{lemma}
\begin{proof}
A \emph{based oriented link} is an oriented link with one distinguished
nonsingular point $p$ on one distinguished component; isotopies carry $p$ and
that component. Fix an oriented knot $K$. For a based link $(J,p)$ define
\[
   F_K(J,p)=P_{K\#_pJ},\qquad G_K(J,p)=P_KP_J.
\]
Here $K\#_pJ$ is made in a ball at $p$, disjoint from the rest of $J$, by
deleting a short oriented arc through $p$, cutting a copy of $K$ at a short
oriented arc, and using the orientation-compatible cross-pairing of incoming
and outgoing endpoints. Two sufficiently small choices of ball and arc are
related by an isotopy supported near $p$, and a based ambient isotopy carries
the ball and insertion with it. Thus $F_K$ is well defined on based isotopy
classes. Both $F_K$ and $G_K$ satisfy the HOMFLY--PT skein relation in every
disc disjoint from $p$ (Axiom~\ref{ax:homfly}): the insertion at $p$ is outside
the disc and therefore commutes with the local switch and oriented smoothing.

We use the following relative uniqueness argument, printed because unbased
uniqueness alone does not apply to a marked operation. Suppose two $R$-valued
invariants of based oriented links satisfy that skein relation away from the
mark and agree on every based unlink. Draw a based link diagram. Using $p$ as
the first traversal basepoint, choose auxiliary basepoints on the other
components and order those components after the marked one. Traverse the
oriented components in that order. Call a crossing \emph{bad} when its first
visit is on the underpassing branch. We prove equality by lexicographic
induction on
\[
   (\text{number of crossings},\ \text{number of bad crossings}).
\]
If a bad crossing is present, switch it. Its first visit becomes over, so the
crossing count is unchanged and the bad count drops by one; no other crossing
or first-visit order changes. The oriented smoothing has one fewer crossing.
Its disc misses $p$, so the marked point survives and the unique resulting
component containing it is distinguished; give all remaining components fresh
auxiliary basepoints and an order. For either invariant $Q$, if the original
crossing is $L_+$, solve the skein relation as
\[
   Q(L_+)=a^{-2}Q(L_-)+a^{-1}zQ(L_0);
\]
if it is $L_-$, solve it as
\[
   Q(L_-)=a^2Q(L_+)-azQ(L_0).
\]
Only the units $a^{\pm1}$ are used. The induction hypotheses for the switched
and smoothed diagrams give equality at the original diagram.

If there is no bad crossing, the ordered based diagram is descending. Keep a
small arc at the current component's basepoint fixed. Starting just after it
and following the orientation, pull each successive already-traversed arc
upward and across the still-untraversed diagram. At every crossing that moving
arc is the overpassing branch, so the pull crosses no strand. Successively
remove all self-crossings of that component; because its visits precede those
of every later component, it also passes over every later component and can be
slid into a disjoint ball as a round circle. For the first component the fixed
arc contains $p$, so this is rel the mark. Repeat in the chosen component
order. The result is the based unlink with the same number of components.
Thus agreement on based unlinks closes the induction.

Apply relative uniqueness first to $F_K,G_K$. On the based $r$-component
unlink $U_r$, with $U_0$ denoting no extra split component,
Definition~\ref{def:homfly} gives
\[
 F_K(U_r,p)=P_{K\sqcup U_{r-1}}=\delta^{r-1}P_K
 =P_KP_{U_r}=G_K(U_r,p).
\]
Hence $P_{K\#_pJ}=P_KP_J$. For the split operation put $K$ in a ball
disjoint from $J$ and define
\[
 F_K^{\sqcup}(J,p)=P_{K\sqcup J},\qquad
 G_K^{\sqcup}(J,p)=\delta P_KP_J.
\]
These again satisfy the skein relation away from $p$, and repeated use of the
split-circle identity in Definition~\ref{def:homfly} gives, on $U_r$,
\[
 F_K^{\sqcup}(U_r,p)=P_{K\sqcup U_r}=\delta^rP_K
 =\delta P_KP_{U_r}=G_K^{\sqcup}(U_r,p).
\]
Relative uniqueness proves $P_{K\sqcup J}=\delta P_KP_J$ for every based
oriented link $J$; this is the split identity of~(i). Clause~(iii) follows by
induction on $n$: the case $n=1$ is trivial, and
$K_1\sqcup(K_2\sqcup\dots\sqcup K_n)$ applies~(i)'s split identity with knot
$K_1$ and link $K_2\sqcup\dots\sqcup K_n$.

For~(ii), put $h(E)=[z^{-1}]P_E$. At a mixed crossing of a two-component
diagram, the oriented smoothing $E_0$ is a knot. Extracting the $z^{-1}$
coefficient from the skein relation gives
\[
 a h(E_+)-a^{-1}h(E_-)=[z^{-2}]P_{E_0}=0.
\]
The last equality uses the full retained clause
$P_{E_0}\in\mathbb Z[a^{\pm1},z^2]$: this is a polynomial ring in $z^2$, so
its $z$-powers are nonnegative as well as even. Thus
$h(E_+)=a^{-2}h(E_-)$. With the convention
$\lambda_+=\lambda_-+1$ at a positive versus negative mixed crossing,
$a^{2\lambda(E)}h(E)$ is invariant under every mixed crossing switch.
No self-crossing is switched: its oriented smoothing has three components,
not a knot, so the displayed vanishing is unavailable.

Switch precisely the mixed crossings at which component $D_1$ passes under
$D_2$. In the resulting diagram $D^\uparrow$, $D_1$ passes over $D_2$ at every
mixed projected intersection. Realize its self-crossings by strict local height
differences. Translate the whole first component upward and the whole second
component downward. Every mixed height gap strictly increases, the
self-crossing gaps are unchanged, and no other projected pair can collide.
This is an isotopy through embeddings; for a large translation a horizontal
plane separates the two components. Their knot types are still $D_1,D_2$ and
the endpoint is split, so its linking number is zero. Mixed-switch invariance
and~(i) now give
\[
\begin{aligned}
 h(D)
 &=a^{-2\lambda}h(D^\uparrow)
  =a^{-2\lambda}[z^{-1}]
       \bigl(\delta P_{D_1}P_{D_2}\bigr)\\
 &=a^{-2\lambda}(a-a^{-1})
       [z^0]\bigl(P_{D_1}P_{D_2}\bigr),
\end{aligned}
\]
which is~(ii).
\end{proof}

\begin{corollary}[interlacement connected sums and the grouped carrier]\label{cor:groupedknot}
\begin{enumerate}
\item[(A)] Let $S\in\Ind(G_P)$, let $L$ be a carrier of $S$ bearing at least one residual
piece, and let
\[
   W=\bigcup\{\,H: H\text{ a residual piece of }S\text{ carried by }L\,\}
\]
be the \emph{union of their label sets} --- a set of crossing labels, not a set
of pieces. By Lemma~\ref{lem:piececurve} Step 2, $W$ is exactly the set of
self-crossings of $L$, so the diagram $D(W)$ obtained from $L$ by retaining
exactly the labels of $W$, each resolved by the divide convention, and erasing
all others, is $L$ itself carrying all of its own crossings; no realization
argument is needed. Let $W=W_1\sqcup\dots\sqcup W_k$, $k\geq1$, be the partition
of $W$ into the connected components of the interlacement graph induced on $W$;
these are the label sets of the residual pieces. The components can be indexed
and equipped with a fixed ordered binary parenthesization $\mathcal T$ such that
\[
D(W)\;\cong\;\#_{\mathcal T}
   \bigl(D(W_1),\dots,D(W_k)\bigr),
\qquad
P_{D(W)}=\prod_{i=1}^{k}P_{D(W_i)} .
\]
Here $\#_{\mathcal T}$ is evaluated recursively: a leaf has value $D(W_i)$
and an internal node has the connected sum of its two ordered child values
with the displayed parenthesization. In particular, when $k=1$,
$\#_{\mathcal T}(D(W_1))=D(W_1)$.
\item[(B)] Let $S\in\Ind(G_P)$ and let $L$ be a carrier of $S$ bearing the residual pieces
$H_1,\dots,H_k$ with $k\geq1$. (If $L$ bears no piece the assertion is
Theorem~\ref{thm:carrierfloor}(D): $P_{S,L}=1$, $w_{S,L}=0$, and there is nothing to
decompose. Clause~(A) is stated for a nonempty label set and
is not applied in that case.) Then the diagram obtained from $L$ by retaining exactly the
crossings of $H_1\cup\dots\cup H_k$ is a knot diagram whose HOMFLY polynomial is
$P_{S,L}=\prod_iP_{H_i}$, whose writhe is $w_{S,L}=\sum_i|H_i|$, and whose
underlying plane curve is $L$, which has no triple points.
\end{enumerate}
\end{corollary}
\begin{proof}
\emph{Clause (A).}
Read the carrier as a cyclic word in the labels of $W$, each label occurring
twice. Call a \emph{gap} a position between two consecutive visits.

\emph{Gap lemma.} Fix a connected component $W_i$. A gap has no intrinsic pair
of sides, so the coordinate is taken per chord: for a gap $g$ and a chord
$x\in W_i$, record which of the two components of the traversal circle minus the
two visits of $x$ contains $g$. Ranging over $x\in W_i$ this is a binary vector
indexed by the chords of $W_i$, the \emph{side vector} of $g$. Distinct gaps of $W_i$ have
distinct side vectors. Suppose two distinct gaps had the same vector. The
oriented interval between the two cuts then contains either both or neither
endpoint of every chord of $W_i$; since the gaps are distinct, both the
interval and its complement contain visits of $W_i$. So the labels of $W_i$
split into two nonempty sets, one with both endpoints inside and one with both
endpoints outside, and no chord of one set interlaces a chord of the other.
That disconnects $W_i$, contrary to hypothesis.

\emph{Every other component sits inside one gap of $W_i$.} A chord outside
$W_i$ fails to interlace every chord of $W_i$ exactly when its two endpoint
gaps have the same side vector, hence, by the gap lemma, exactly when both its
endpoints lie in one and the same gap of $W_i$. If some component $W_j$,
$j\neq i$, occupied two different gaps of $W_i$, its chords in the two gaps
could not interlace each other, contradicting connectedness of $W_j$.

\emph{Closed-block tree.} Apply the preceding two paragraphs inside any
cyclic node word containing at least two of the components. They give a
nonempty proper label-closed linear block $\alpha$; its complementary linear
word $\beta$ is label-closed as well. Record the two closed cyclic child words
in the order $(\alpha,\beta)$. In each child, identifying its two cut ends
creates one distinguished \emph{join gap}. Both children have fewer labels,
so iteration gives a finite ordered binary tree $\mathcal T$ whose leaves are
exactly the $W_i$. Index the leaves in its order. Every non-root node also
remembers the one gap in its cyclic record prescribed when its parent was cut;
that output gap is distinct data from the two child join gaps used when the
node itself is assembled.

We induct upward with the following invariant. A node $V$ has a connected
generic oriented one-circle diagram $C_V$ and an orientation-preserving record
isomorphism from its record to the cyclic record of $V$; if $V$ is not the
root, it also has a clean open subarc in the exact output gap prescribed by its
parent.

For a leaf $V=W_i$, take the actual curve $C_{W_i}$ constructed by
Lemma~\ref{lem:piececurve}, with the inherited divide crossing data. It has
exactly the distinct inherited transverse double points of $W_i$, its
restricted cyclic word, and no triple point: the construction only smooths in
pairwise disjoint crossing discs and creates no intersection.
Lemma~\ref{lem:pieceintrinsic} supplies the record comparison with the carrier
factor $D(W_i)$. If the leaf is not the root, the extension clause in
Definition~\ref{def:record} carries its parent-prescribed output gap to a
complementary arc of its traversal circle. Choose an interior regular point
there and a sufficiently small subarc about it; a small disc meets the curve
in that arc alone, so it is clean. If the leaf is the root, the invariant
requires no output gap.

For an internal node with ordered children $A,B$, use their clean join arcs.
Take $C_A$ as the based host $J$, with $p$ in its join arc, and take $C_B$ as
$K$. At an interior point of $K$'s join arc choose an
orientation-preserving smooth chart of the sphere sending that point to
infinity. A sufficiently small spherical cap about infinity meets $K$ in one
embedded arc, so its complement is a compact long $K$-tangle with two oriented
ends. Shrink that tangle into the clean host disc, delete the short host arc
through $p$, place the tangle ends next to their prescribed host ends, and
cross-pair outgoing $A$ to incoming $B$ and outgoing $B$ to incoming $A$ by
two disjoint short arcs. Round the two joins inside disjoint smaller discs.
This is literally the orientation-compatible local insertion $K\#_pJ$ of
Lemma~\ref{lem:homflyrows}(i), not an unnamed planar splice.

The insertion disc contains no other point of $J$ and the inserted tangle and
connectors lie inside it. Thus it creates no crossing. The sphere chart is
orientation preserving, so all child crossing signs are preserved, and the
over/under designations are transported. Starting immediately before $A$, the
one traversal circle reads the linear word $\alpha$ and then $\beta$.
Pairings stay within a child. These are precisely conditions (a)--(d) of
Definition~\ref{def:record}. The child crossings remain transverse and
distinct, the connectors are disjoint, and the local rounding creates no
intersection, so the result is a generic one-circle diagram with no triple
point.

It remains to mark the node's output gap when the node is not the root. If that
gap is not one of the two split boundaries, it survives in the appropriate
child and is carried across by the record map. If it is a split boundary, use
the corresponding new connector subarc. In either case a smaller disc meeting
only that arc makes it clean. This proves the invariant.

Finally the actual diagram $D(W)$ has underlying curve $L$. The diagrammatic
parent has isolated transverse double points with no two sharing an image, so
the smoothing sites admit pairwise disjoint discs, each meeting the curve in
the two crossing arcs alone. Each local smoothing replaces them by two
noncrossing arcs and outside those discs the curve is unchanged. Hence $L$ is
a connected generic one-circle curve with no triple point. The constructed
root has the same properties by the invariant, and its record is isomorphic to
the record of $D(W)$: the traversal order is the tree's ordered concatenation,
labels and pairings never cross between children, and every over/under role and
sign is inherited. Axiom~\ref{ax:gausscode} identifies the two oriented links.
At the leaves Lemma~\ref{lem:pieceintrinsic} identifies $C_{W_i}$ with
$D(W_i)$. Thus the root is exactly the displayed
$\mathcal T$-parenthesized connected sum, without an associativity assumption.

The polynomial identity follows by applying
Lemma~\ref{lem:homflyrows}(i) at each internal node of the fixed tree
$\mathcal T$; every leaf is a knot diagram. Thus
$P_{D(W)}=\prod_iP_{D(W_i)}$.

\emph{Clause (B).}
Each $H_i$ has both visits of each of its labels on $L$ by
Lemma~\ref{lem:carriers}(iv), and the $H_i$ are exactly the connected
components of the interlacement graph induced on $H_1\cup\dots\cup H_k$, since
they are components of $G_P[U(S)]$ and interlacement is read from the same
word. Clause~(A) applies. The diagram has one component
because $L$ is one closed curve, the writhe is the crossing count because every
crossing is positive (Definition~\ref{def:piecediagram}), and the underlying
plane curve is $L$ because erasing a double point omits it from the record and
performs no operation on the curve (Definition~\ref{def:piecediagram}).

The generic parent polygon is diagrammatic by
Definition~\ref{def:diagrammatic}, so it has no triple points.  The parent's
self-intersections are isolated and no two share an image
(Definition~\ref{def:diagrammatic}), so they admit pairwise disjoint discs,
each meeting the curve in the two crossing arcs alone.  Inside each disc the
oriented smoothing of Definition~\ref{def:smoothing} replaces those two arcs
by two arcs that do not cross, and outside the discs the curve is unchanged;
so no intersection is created and $L$ has no triple points either.

The identification of the factors is Lemma~\ref{lem:pieceintrinsic}: the factor
$D(H_i)$ produced by clause~(A) is the diagram of $L$ with
only the crossings of $H_i$ retained, and the identity on the visits of $H_i$ is
a record \emph{isomorphism} from that diagram's record to the record of the
intrinsic piece diagram of $H_i$ on the parent polygon --- not an equality, the
two traversal circles being different --- so by Axiom~\ref{ax:gausscode} the two
present the same oriented link, whose polynomial is $P_{H_i}$ and whose writhe
is $|H_i|$. So the product in the display is a
product of the $P_{H_i}$ and not of some other family. An earlier revision
asserted this identification in one sentence, on the ground that both
constructions ``retain the crossings of $H_i$''; the peer refuted that ground
--- smoothing the support changes the traversal successor, so containment of
both visits does not identify the cyclic records --- and
Lemma~\ref{lem:pieceintrinsic} now supplies the three components of the record
separately.
\end{proof}

\begin{remark}[why clause (A) is not stated for an arbitrary label set]
\label{rem:connectedsumscope}
An earlier revision quantified clause~(A) over any set $W$ of labels with both
visits on the carrier. That is not well posed. Restricting a realizable Gauss
word to an arbitrary label set can produce a word no plane curve realizes: the
peer's witness restricts $123123$ to two of its three crossings and obtains
$2323$, in which each chord interlaces exactly one other, while every chord of a
realizable word interlaces an even number. For such a $W$ the diagram $D(W)$
does not exist and neither side of the display means anything.

For the full self-crossing set $W$ of a carrier --- the hypothesis of clause~(A) ---
no realization theorem is needed at all: by Lemma~\ref{lem:piececurve} Step 2 the
undominated crossings carried by $L$ are exactly the self-crossings of $L$, so
$D(W)$ is the carrier itself. An intermediate revision instead allowed $W$ to be
any nonempty \emph{union} of residual pieces and cited
Lemma~\ref{lem:piececurve} to realize it; the peer refused that citation and was
right --- that lemma realizes one \emph{connected} piece, its Step~5 requiring
connectedness, and a disconnected union is outside it. The citation is
withdrawn; what replaces it is the observation that the set
$W=H_1\cup\dots\cup H_k$ supplied by clause~(B) is the whole self-crossing
set of an actual curve.

The witness is excluded for a separate and still valid reason: two interlacing
crossings that interlace nothing else cannot be a residual piece, since by
Corollary~\ref{cor:evendominated} an undominated crossing meets an even number
of dominated crossings, so its degree inside its own piece is even, and in a
two-element piece it would be one.
\end{remark}

\begin{remark}[a generalization that was tried and withdrawn]\label{rem:groupedgeneral}
An intermediate revision quantified clause~(B) over an arbitrary closed
curve rather than a carrier, so that it could be applied to a sub-loop cut out
of a carrier at one of its crossings, and discharged the resulting comparison
of diagrams by asserting that the complementary excursion, carrying only erased
crossings, could be contracted. The peer rejected that twice and was right
both times: the general statement does not bind the parent curve or the map to
the intrinsic piece diagrams of Definition~\ref{def:piecediagram}, and an
embedded arc whose interior is disjoint from a diagram need not be
boundary-parallel, so ``separated in space'' does not produce the disc the
contraction needs. Clause~(B) is therefore back in its carrier form, and
Theorem~\ref{thm:s7universal}(D)(i), which wanted the general form,
is reproved below by passing to a larger support --- whose carriers \emph{are}
the two loops --- so that only carrier instances of clause~(B) are used.
\end{remark}

\begin{remark}[why this is needed]\label{rem:whygrouped}
Theorem~\ref{thm:carrierfloor}(C) bounds $\min\deg_a$ of one knot diagram in terms
of \emph{its} writhe and rotation. Corollary~\ref{cor:groupedknot}(B) says the
grouped object $P_{S,L}$, with its grouped writhe $w_{S,L}$ and the carrier's
own rotation $R(L)$, is exactly such a diagram. Without it the floor would have
to be proved for each piece separately and then convolved, which is false in
general: the sum of the pieces' floors is not the floor of the product.
\end{remark}

\subsection{The contact-turn geometry}

\begin{lemma}[the contact turns, both branches]\label{lem:contactturnboth}
Under Convention~\ref{conv:s7standing}, let a simple vertex-on-edge event of
bigon type have moving vertex $M$, remote edge $(a,b)$, and neighbours
$p=p_{M-1}$ and $q=p_{M+1}$.
Let $\epsilon=1$ if the two newborn crossings
(Convention~\ref{conv:s7standing}) interlace, and $\epsilon=0$
otherwise. Let $s_{\mathrm{low}}$ denote the sign \eqref{eq:S6} evaluated in the low
(zero-newborn) chamber. First take an orientation-preserving affine
chart sending the contact to the origin and the direction of $a\to b$ to
$+x$, and write $p=(u,h)$ and $q=(v,k)$. Then $hk>0$. If $h,k<0$,
compose with reflection in the $x$-axis; thereafter work in the resulting
upper picture, where $h,k>0$ and $0<\arg p,\arg q<\pi$. This reflection
reverses all three turns below and $s_{\mathrm{low}}$, so the sign comparisons are
independent of which side was reflected. In the upper picture:
\begin{enumerate}
\item[(i)] $\epsilon=0$ exactly when $\arg p<\arg q$, and $\epsilon=1$
exactly when $\arg q<\arg p$.
\item[(ii)] The two incident directions $-p,q$ are nonzero and transverse to
the remote direction. The contact turns of $\lambda_1,\lambda_2$ are,
respectively, $\arg q$ and $\pi-\arg p$; both lie in $(0,\pi)$ and
have the sign of $s_{\mathrm{low}}$.
\item[(iii)] The full turn at $M$ is the principal representative of
$\arg q-\arg p+\pi$. It is
$\arg q-\arg p-\pi\in(-\pi,0)$ when $\epsilon=0$, with sign
$-s_{\mathrm{low}}$, and $\arg q-\arg p+\pi\in(0,\pi)$ when
$\epsilon=1$, with sign $s_{\mathrm{low}}$.
\item[(iv)] Each half polygon inherits exactly one of these contact turns:
$\lambda_1$ closes along the remote direction and leaves along $q$, while
$\lambda_2$ enters along $-p$ and leaves along the remote direction.
Consequently, when $\epsilon=0$, any half contact carrier whose remaining
turns are inherited from a uniformly turning full contact carrier has exactly
one turn opposite to the full turn, and that dissenting turn has magnitude
strictly below $\pi$.
\end{enumerate}
\end{lemma}
\begin{proof}
For a sufficiently small displacement into the two-newborn chamber, normalize the
remote line to $y=0$ and write
$M_t=(m_t,-c_t)$ with $c_t>0$, and the two neighbours as
$p_t=(u_t,h_t)$ and $q_t=(v_t,k_t)$, tending to $p=(u,h)$ and $q=(v,k)$,
respectively. The incident-edge intersections with the remote line have abscissae
$x_p=m_t+c_t(u_t-m_t)/(h_t+c_t)$ and
$x_q=m_t+c_t(v_t-m_t)/(k_t+c_t)$, so
$(x_p-x_q)/c_t\longrightarrow u/h-v/k=(uk-hv)/(hk)$.
The vertex turn is nonzero, hence $uk-hv\neq0$, so this sign is fixed for small
$t$. Along the polygon the incoming crossing precedes the outgoing crossing;
along the oriented remote edge their order is increasing abscissa. They are
therefore noninterlaced exactly when $x_p>x_q$, and interlaced when $x_p<x_q$.
In the upper half-plane $x_p>x_q$ is exactly $\arg p<\arg q$. This proves (i).

The turn at $M$ is from incoming direction $-p$ to outgoing direction $q$,
hence is represented by $\arg q-\arg p+\pi$. If $\epsilon=0$, clause (i)
gives $\arg p<\arg q$, so its principal value is
$\arg q-\arg p-\pi\in(-\pi,0)$. If $\epsilon=1$, clause (i) gives
$\arg q<\arg p$, so the displayed representative lies in $(0,\pi)$.
In the upper low chamber $M$ is on the neighbours' side of the remote line,
so $s_{\mathrm{low}}=+1$ by \eqref{eq:S6}; this proves (iii),
and reflection negates both sides of every claimed sign comparison.

By Definition~\ref{def:halves}, $\lambda_1$ closes in direction $+x$ and
then leaves along $q$, while $\lambda_2$ arrives along $-p$ and then leaves
in direction $+x$. Their turns are $\arg q$ and $\pi-\arg p$, proving (ii)
and the inheritance assertion in (iv). In the noninterlacing case (iii)
makes the full contact turn opposite to both half contact turns. If the full
carrier is uniform (Definition~\ref{def:wind}), every remaining turn inherited by
either half has the full turn's sign; its one contact turn has the opposite
sign and magnitude in $(0,\pi)$. This proves (iv).
\end{proof}

\subsection{The two sectors of the bigon jump}\label{sec:sectors}

Fix a simple vertex-on-edge event of bigon type: the two neighbours of the
moving vertex $M$ lie strictly on the same side of the line of the remote edge
$(a,b)$, so that crossing the wall changes the crossing number by two. Orient
the parameter from the chamber with no newborn crossing to the chamber with
two; write $P_0$ for the first and $P_2$ for the second, and let
$\delta=+1$ in that direction and $-1$ in the reverse. By
Remark~\ref{rem:s7wellposed} nothing depends on the choice.

\begin{definition}[newborns, sectors, and the interlacement bit]\label{def:sectors}
Under Convention~\ref{conv:s7standing} --- in particular with both halves
generic, without which the products $X_1(\lambda_1)X_1(\lambda_2)$ compared
below are not defined --- the newborn crossings $x,y$ of
Convention~\ref{conv:s7standing} are realized in $P_2$. The directed jump is
$F=\delta\bigl(X_1(P_2)-X_1(P_0)\bigr)$. Split the supports of $P_2$ by how
many newborns they contain and write
\[
   F=R+\mathcal B,
\]
where $\mathcal B$ is $\delta$ times the total contribution of the supports of
$P_2$ containing \emph{both} newborns and $R$ is everything else, the
\emph{returned} sector. Let $\epsilon=1$ if $x$ and $y$ interlace in $P_2$ and
$\epsilon=0$ otherwise, and let $J=s\,X_1(\lambda_1)X_1(\lambda_2)$ be the
required jump, with $s$ the sign \eqref{eq:S6} taken in the initial chamber.
\end{definition}

\begin{lemma}[the contact word]\label{lem:contactword}
The cyclic Gauss word of $P_2$, rotated to begin at the first visit of the
newborn $x$ --- the rotation is named here because consumers speak of the visit
that \emph{opens} the word, which a cyclic word does not have until one is
fixed --- reads
\begin{equation}
   W_2=\begin{cases}
     x_0\,y_0\,A\,y_1\,x_1\,B,&\epsilon=0,\\
     x_0\,y_0\,A\,x_1\,y_1\,B,&\epsilon=1,
   \end{cases}
   \qquad\text{and in both cases}\qquad
   W_0=A\,B
   \label{eq:contactword}
\end{equation}
where $x_0,x_1$ are the two visits of $x$, $y_0,y_1$ those of $y$, $A$ is the
subword of old labels traversed from $M+1$ to $a$ together with the part of the
remote edge from $a$ to \emph{the first newborn visit on that edge} --- which is
$y$ when $\epsilon=0$ and $x$ when $\epsilon=1$, as the two displayed words show
--- and $B$ is the subword traversed from $b$ to $M-1$. An earlier revision said
``to $y$'' in both branches, which is false in the interlacing one; the proof
always used the order-neutral phrase. Moreover an old label $q$ has
\begin{enumerate}
\item[(i)] both visits in $A$, or both in $B$, or exactly one in each;
\item[(ii)] and it is adjacent in $G_{P_2}$ to $x$, or to $y$, exactly when it
has one visit in each --- so the labels adjacent to $x$ and those adjacent to
$y$ coincide, and form the \emph{common neighbourhood} $\mathcal N$.
\end{enumerate}
\end{lemma}
\begin{proof}
In $P_2$ the two edges incident to $M$ each meet the remote edge once, at $x$
and at $y$ respectively; the contact disc contains no other vertex or crossing,
because a further incidence at the event would be a vanishing guarded predicate
outside the zero set that Convention~\ref{conv:events} allows, which is exactly
$\{\mathrm{G2}_{(a,b),M},\ \mathrm{G4}\langle (a,b);e_{M-1},e_M\rangle\}$ ---
the member-valued accessor, a zero set holding members
(Definition~\ref{def:guarded}) --- and contains no other incidence. Traversing from just before $M$: the
edge $M-1\to M$ carries the visit $x_0$, then comes the vertex $M$, then the
edge $M\to M+1$ carries $y_0$; no crossing lies between them. Continuing, the
traversal runs $M+1,\dots,a$ and then along the remote edge, where it meets the
two newborns again, in one of the two orders, before continuing $b,\dots,M-1$.
Writing $A$ for the stretch from just after $y_0$ to just before the first
remote-edge newborn visit and $B$ for the stretch from just after the second to
just before $x_0$, the word is $x_0y_0Ay_1x_1B$ or $x_0y_0Ax_1y_1B$ according
to that order. In the first case the interval from $x_0$ to $x_1$ is
$y_0Ay_1$, which contains both visits of $y$, so $x$ and $y$ do not interlace;
in the second it is $y_0A$, which contains exactly one, so they do. The two
cases are therefore $\epsilon=0$ and $\epsilon=1$, which is
\eqref{eq:contactword}. Deleting
the four newborn visits leaves $AB$, and $P_0$ has exactly the old crossings of
$P_2$ with the same order, since the two chambers' crossing sets differ by
exactly the newborn pair (the argument of Lemma~\ref{lem:cuspcases}(iv), with the
vertex-on-edge zero in place of the cusp one). That is $W_0=AB$.

For (i) and (ii), take $\epsilon=0$ first: the interval from $x_0$ to $x_1$ is
$y_0Ay_1$ and the interval from $y_0$ to $y_1$ is $A$. An old label with both
visits in $A$ has two visits in each interval, so it interlaces neither
newborn; with both visits in $B$ it has none in either, so again neither; with
one in $A$ and one in $B$ it has exactly one in $y_0Ay_1$ and exactly one in
$A$, so it interlaces both. For $\epsilon=1$ the two intervals are $y_0A$ and
$Ax_1$, and the same three cases give two, zero and one visits respectively, so
the conclusion is identical.
\end{proof}

\begin{definition}[eligible]\label{def:eligible}
An old support $T\in\Ind(G_{P_0})$ is \emph{eligible} when $T\cap\mathcal
N=\varnothing$, i.e.\ when no element of $T$ is adjacent to a newborn.
\end{definition}

\begin{remark}[why eligibility is stated this way and not the other]
\label{rem:eligibility}
The tempting definition, and the one an earlier revision of this section used,
is "$T\cup\{x,y\}$ is independent". In the branch $\epsilon=0$ the two agree:
$x$ and $y$ are then non-adjacent to each other, and $T$ is independent, so
$T\cup\{x,y\}$ is independent exactly when $T$ misses the neighbours of both.
In the branch $\epsilon=1$ they do not agree at all --- there $x\sim y$, so
$T\cup\{x,y\}$ is \emph{never} independent and that definition has empty
domain. Every statement below about the returned sector would then be
vacuous at $\epsilon=1$ and would prove nothing about the interlacing bigon.
Definition~\ref{def:eligible} is the condition the arguments actually use: it
is what makes the old labels split between the two intervals, and it is
non-vacuous in both branches.
\end{remark}

\begin{lemma}[the eligible old supports are an ordered product]\label{lem:orderedproduct}
Here eligibility is as in Definition~\ref{def:eligible}.
Let $T$ be eligible. Then, writing $V_A$ and $V_B$ for the old labels with both
visits in $A$ and in $B$,
\[
   G_{P_2}\bigl[V_A\sqcup V_B\bigr]=G_A\sqcup G_B
\]
is a disjoint union with no edge between the parts, and
$T\mapsto(T\cap V_A,\;T\cap V_B)$ is a bijection from the eligible supports
onto $\Ind(G_A)\times\Ind(G_B)$.
\end{lemma}
\begin{proof}
An eligible $T$ contains no element of $\mathcal N$ by
Definition~\ref{def:eligible}, and by Lemma~\ref{lem:contactword}(ii)
$\mathcal N$ is exactly the set of old labels with one visit in each interval,
so by Lemma~\ref{lem:contactword}(i) every element of $T$ lies in $V_A$ or in
$V_B$. A chord of $V_A$ has both endpoints in the cyclic interval $A$ and a
chord of $V_B$ both endpoints in the disjoint interval $B$, so their endpoints
cannot alternate: there is no edge between the parts. Independent sets of a
disjoint union are exactly pairs of independent sets of the parts, and the map
displayed is the inverse of $(T_A,T_B)\mapsto T_A\sqcup T_B$. The order of the
pair is the order of the halves of Definition~\ref{def:halves}, not a quotient:
$A$ is the word of $\lambda_1$ and $B$ that of $\lambda_2$, because $\lambda_1$
runs $M\to M+1\to\dots\to a$ and closes along the remote edge, which is the
stretch $A$, and $\lambda_2$ opens along the remote edge and runs
$b\to\dots\to M-1$, which is $B$.
\end{proof}

\subsection{The noninterlacing both-newborn support calculation}\label{sec:niboth}

\begin{proposition}[the noninterlacing carrier factorization and sector target]\label{prop:sectortarget}
For clauses~(A)--(C), let $\epsilon=0$, let $T$ be an eligible old support
with ordered parts $T_A,T_B$ as in Lemma~\ref{lem:orderedproduct}, and put
$S=T\cup\{x,y\}$. Write $\mathcal N$ for the common neighbourhood of
Lemma~\ref{lem:contactword}, and for a word $W$ write
$\mathrm{Del}_{\mathcal N}(W)$ for $W$ with every visit of a label of
$\mathcal N$ deleted; put $A^{\circ}=\mathrm{Del}_{\mathcal N}(A)$ and
$B^{\circ}=\mathrm{Del}_{\mathcal N}(B)$. Let $s_{\mathrm{low}}$ be the sign
\eqref{eq:S6} computed in the low chamber $P_0$. For clauses~(D)--(F), let
$F,R,\mathcal B,\epsilon,J$ be as in Definition~\ref{def:sectors}. Then:
\begin{enumerate}
\item[(A)] The carriers of $S$ in $P_2$ are, canonically and without repetition, the
carriers of $T_A$ in $\lambda_1$, the carriers of $T_B$ in $\lambda_2$, and one
further carrier $U_{\mathrm{loc}}$ lying in the contact disc.
\item[(B)] Under the correspondence of clause~(A), each carrier of $S$
other than $U_{\mathrm{loc}}$ has the same residual pieces, the same signed piece words,
the same piece polynomials, the same grouped polynomial $P_{S,L}$ and the same
grouped writhe $w_{S,L}$ as its mate in the half. The local carrier bears no
piece: $P_{S,U_{\mathrm{loc}}}=1$ and $w_{S,U_{\mathrm{loc}}}=0$.
\item[(C)] Here $\rot$ is as in Definition~\ref{def:rot}.
Under the correspondence of clause~(A), each carrier of $S$ other than $U_{\mathrm{loc}}$ has the same cyclic
corner list with the same turn signs, hence the same weight, signed rotation,
$R$, and turn-uniformity as its half mate. The local carrier has
\[
   P_{S,U_{\mathrm{loc}}}=1,\quad w_{S,U_{\mathrm{loc}}}=0,\quad \rot(U_{\mathrm{loc}})=-s_{\mathrm{low}},\quad R(U_{\mathrm{loc}})=1,\quad
   \Omega_1(S,U_{\mathrm{loc}})=1,\quad \mathrm{wt}(U_{\mathrm{loc}})=s_{\mathrm{low}} .
\]
\item[(D)] if $\epsilon=1$, then $\mathcal B=0$;
\item[(E)] if $\epsilon=0$, then $\mathcal B=J$;
\item[(F)] in both cases $\mathcal B=(1-\epsilon)J$, and therefore
\[
   F=J\quad\Longleftrightarrow\quad R=\epsilon J .
\]
\end{enumerate}
In particular the bigon branch of (S7) holds if and only if the returned sector
is $J$ when the newborns interlace and $0$ when they do not.
\end{proposition}
\begin{proof}
\emph{Clause (A).}
Smoothing reconnects the traversal by the successor rule: at a smoothed
crossing with visits $c_0,c_1$, the arc arriving at $c_0$ continues with the
arc leaving $c_1$, and conversely. Applied to \eqref{eq:contactword} at the two
newborns this gives
\[
   \sigma(x_0)=B_0,\qquad
   \sigma(x_1)=y_0,\qquad
   \sigma(y_0)=x_1,\qquad
   \sigma(y_1)=A_0,
\]
where $A_0$ and $B_0$ are the first positions of $A$ and of $B$ --- and, when
one of those words is empty, the adjacent marker, namely $y_1$ for an empty
$A$ and $x_0$ for an empty $B$. The middle two are a $2$-cycle
$\{y_0,x_1\}$: that is the local carrier $U_{\mathrm{loc}}$, the closed curve made of the
two smoothing arcs at $x$ and $y$ together with the stretch of the remote edge
between them and the two incident stretches through $M$. The other two
successors leave $x_0$ into $B$ and $y_1$ into $A$, and by
Lemma~\ref{lem:orderedproduct} the selected old labels of $A$ and of $B$ are
$T_A$ and $T_B$, whose own smoothing is internal to their intervals. So the
remaining cycles are one family confined to $A$ and one confined to $B$.

To match them with the halves, delete from each the designated newborn boundary
marker and every visit of a label of $\mathcal N$. Such a label crosses between
the two halves, so it appears in neither half's self-crossing word; and it is
not smoothed, since an eligible support misses $\mathcal N$, so its visits are
through-points and not corners --- deleting them removes no corner. A step that
lands on the deleted boundary marker becomes the cyclic wraparound of the
cleaned word. The projected cycles are therefore exactly the smoothing cycles
of $T_A$ on $A^{\circ}$ and of $T_B$ on $B^{\circ}$, which are the carrier
books of $\lambda_1$ and $\lambda_2$ by Lemma~\ref{lem:carriers} applied inside
each half. If a cleaned word is empty its single boundary cycle projects to the
one crossing-free carrier of that half, which is that half itself.


\emph{Clause (B).}
A label of $\mathcal N$ is adjacent to both newborns
(Lemma~\ref{lem:contactword}(ii)), so in $P_2$ it lies in
$N_{G_{P_2}}(S)$ and is absent from $U(S)$; geometrically it is a crossing between
the two halves and appears in neither half's diagram. Every old label outside
$\mathcal N$ has both visits in one interval, so its interlacements with other
such labels and its cyclic visit order are those of that half. The over/under
role of each persistent old crossing is the sign of its $\mathrm{G5}$
determinant, which is unchanged across the event: the crossing persists, so that
member is \emph{active} and hence relevant, and a relevant member outside the
zero set keeps its sign by Lemma~\ref{lem:guardconst}. An earlier revision
called $\mathrm{G5}$ unconditional, which it is not. By Lemma~\ref{lem:carriers}(iv) each
residual component sits on the unique carrier holding all its visits, and
clause~(A) identifies that carrier with its half mate. Hence the signed residual words agree; equal records give equal links by
Axiom~\ref{ax:gausscode}, and hence equal polynomials, and the writhes agree
because they are crossing counts. An earlier revision cited
Axiom~\ref{ax:homfly} for the passage from equal records to equal polynomials,
which is the axiom that computes a polynomial, not the one that identifies the
link a record presents. The cycle
$\{y_0,x_1\}$ contains only newborn visits, so it carries no residual label,
giving the last sentence.


\emph{Clause (C).}
Choose a contact disc containing no old vertex, no old crossing and no
smoothing site, which exists because the contact is isolated. Inside it, deform
each smoothed arc of $P_2$ to the corresponding closure of the half, keeping
the disc boundary fixed. Along the deformation every corner of the carrier is a
polygon vertex or a smoothing site of a crossing of $S$, and no corner
determinant vanishes: the vertex turns are guarded by (G1), the smoothing-site
turns by the active (G5) of their crossings, and the two directions incident to
the contact are nonzero and transverse to the remote direction by
Lemma~\ref{lem:contactturnboth}(ii). Ordinary crossing visits are not corners, so
nothing inside the disc passes through a degenerate turn. The cyclic corner list is fixed along the deformation. Every incident
edge is nonzero and no consecutive pair is a negative multiple, because the
corresponding corner determinant was just shown nonzero. Thus the carrier is a
path in its fixed $\mathcal R_c$, and Lemma~\ref{lem:turnlift}(ii) keeps its
signed rotation constant; the corner signs outside the disc and the matched contact turn keep their signs;
their products, hence the weights of Definition~\ref{def:wind}, agree. This gives the first sentence.

The local carrier is the closed triangle with corners at $x$, at $M$ and at
$y$. If $M$ lies to the left of the oriented remote edge in the low chamber,
i.e.\ $s_{\mathrm{low}}=+1$, then in the high chamber the triangle is traversed clockwise,
with three right turns; Lemma~\ref{lem:turnlift}(ii) gives signed rotation
$-1$. If $s_{\mathrm{low}}=-1$ it is traversed counterclockwise with three left turns, and
the same lemma gives signed rotation $+1$. In
both cases $\rot(U_{\mathrm{loc}})=-s_{\mathrm{low}}$ and $R(U_{\mathrm{loc}})=1$. Its weight is $+1$ in the first case
and $(-1)^3=-1$ in the second (Definition~\ref{def:wind}), i.e.\
$\mathrm{wt}(U_{\mathrm{loc}})=s_{\mathrm{low}}$. It bears no piece by clause~(B), so
$P_{S,U_{\mathrm{loc}}}=1$, $w_{S,U_{\mathrm{loc}}}=0$, its slot is $1-0-1=0$ and
$\Omega_1(S,U_{\mathrm{loc}})=[a^0z^0]1=1$.



\emph{Clause (D).} If $\epsilon=1$, then $x$ and $y$ interlace and hence are adjacent in
$G_{P_2}$ (Definition~\ref{def:interlace}). No independent set contains both
(Definition~\ref{def:smoothing}), so the both-newborn sector has no supports
and its empty sum is $\mathcal B=0$.

\emph{Clause (E).} Assume $\epsilon=0$. Fix an arbitrary eligible old support $T$, let
$(T_A,T_B)$ be its ordered parts from Lemma~\ref{lem:orderedproduct}, and put
$S=T\cup\{x,y\}$. Let $s_{\mathrm{low}}$ be the sign \eqref{eq:S6} in $P_0$.
By clauses~(A)--(C), the term of $S$ factors as
\[
\wind(S)\prod_{L}\Omega_1(S,L)
=\Bigl(\mathrm{wt}(U_{\mathrm{loc}})\!\!\prod_{L\neq U_{\mathrm{loc}}}\!\!\mathrm{wt}(L)\Bigr)
 \Bigl(\Omega_1(S,U_{\mathrm{loc}})\!\!\prod_{L\neq U_{\mathrm{loc}}}\!\!\Omega_1(S,L)\Bigr)
=s_{\mathrm{low}}\,t_{\lambda_1}(T_A)t_{\lambda_2}(T_B),
\]
where $t_{\lambda_i}$ denotes the corresponding term of the half state sum,
because every factor other than the local one matches a factor of exactly one
half. The both-newborn supports are exactly the $T\cup\{x,y\}$ with $T$
eligible: independence follows from eligibility and $x\not\sim y$, while the
old part of any both-newborn support misses the common neighbourhood.
Lemma~\ref{lem:orderedproduct} indexes these supports by
$\Ind(G_A)\times\Ind(G_B)$, so finite distributivity gives
\[
 \sum_T s_{\mathrm{low}}t_{\lambda_1}(T_A)t_{\lambda_2}(T_B)
 =s_{\mathrm{low}}\Bigl(\sum_{T_A}t_{\lambda_1}(T_A)\Bigr)
  \Bigl(\sum_{T_B}t_{\lambda_2}(T_B)\Bigr)
 =s_{\mathrm{low}}X_1(\lambda_1)X_1(\lambda_2),
\]
including the empty supports in both complete half state sums. Let
$s_{\mathrm{init}}$ be the sign in the initial chamber used by
Definition~\ref{def:sectors}. The two chambers have opposite side signs; the
initial chamber is $P_0$ for $\delta=+1$ and $P_2$ for $\delta=-1$. Hence
$s_{\mathrm{init}}=\delta s_{\mathrm{low}}$, and the directed sector is
\[
 \mathcal B=\delta s_{\mathrm{low}}X_1(\lambda_1)X_1(\lambda_2)
 =s_{\mathrm{init}}X_1(\lambda_1)X_1(\lambda_2)=J.
\]

\emph{Clause (F).} Clauses (D) and (E), with $\epsilon\in\{0,1\}$, give
$\mathcal B=(1-\epsilon)J$. By Definition~\ref{def:sectors},
$F-J=R+\mathcal B-J=R-\epsilon J$, proving the equivalence.
\end{proof}

\subsection{The returned sector: framing and the spanning recurrence}
\label{sec:spanning}

By Proposition~\ref{prop:sectortarget}(F) the bigon branch of (S7) is exactly the
statement $R=\epsilon J$. This subsection puts the returned sector in a form in
which the interactions between the two contact intervals drop out.

\begin{definition}[the four returned branches]\label{def:branches}
Old labels have the same interlacements in $P_0$ and in $P_2$, since those
depend only on the cyclic order of their four visits, which the newborn visits
do not disturb. So $\Ind(G_{P_0})$ is exactly the set of newborn-free supports
of $P_2$. For $T\in\Ind(G_{P_0})$ write $L(T)$ for its term in $X_1(P_0)$,
$N(T)$ for its term in $X_1(P_2)$, and $X(T)$, $Y(T)$ for the terms of
$T\cup\{x\}$ and $T\cup\{y\}$ in $X_1(P_2)$, each read as $0$ when the set is
not independent. Then
\[
   R=\delta\sum_{T\in\Ind(G_{P_0})}\bigl(N(T)-L(T)+X(T)+Y(T)\bigr).
\]
\end{definition}

\subsection{The selector identities}\label{sec:selectors}



\begin{lemma}[the contact carrier splits]\label{lem:contactsplit}
Here eligibility is as in Definition~\ref{def:eligible}.
Under Convention~\ref{conv:s7standing}, let $T$ be eligible, with parts
$T_A,T_B$. The genericity of the halves is stated here,
where the half objects are first formed, and not only at the endpoint theorems
that consume them: the conclusion speaks of the carriers of $T_A$ in
$\lambda_1$ and of $T_B$ in $\lambda_2$, and a carrier of a support is read off
a polygon's Gauss word and its smoothing sites, which a nongeneric half need
not possess. The hypothesis is (H3) of Convention~\ref{conv:s7standing}, in force for this
section by that convention's declaration; it is named here because this is
where the half objects are first formed, not because naming is the mechanism. In the high-neither state the
carriers of $T$ in $P_2$ are: one distinguished carrier $L^{*}$ containing the
corner $M$ and the remote-edge stretch through the contact, and a family of
others, each of which lies entirely in $A$ or entirely in $B$. Cutting at the
contact carries $L^{*}$ to one carrier $L_1$ of $T_A$ in $\lambda_1$ and one
carrier $L_2$ of $T_B$ in $\lambda_2$, each containing that half's contact
corner, and carries the others bijectively to the remaining carriers of the two
halves.

Five further data are part of the conclusion, because they are used later
and a statement's consumers must find their inputs in its clauses
rather than in a picture:
\begin{enumerate}
\item[(a)] \emph{Corner partition.} Every corner of every carrier of $T$ in
$P_2$ other than the contact corner $M$ of $L^{*}$ --- each being a vertex of
the polygon or a smoothing site, lying strictly inside $A$ or strictly inside
$B$ --- is a corner of exactly one carrier of exactly one half, namely the
half of its interval, under the correspondence above.
\item[(b)] \emph{Ordered direction pairs.} At each such corner the ordered
pair (incoming direction, outgoing direction) read in the half equals the pair
read in $P_2$: the cut alters the curve only at the two contact stretches, and
a corner strictly inside an interval keeps both its incident arcs.
\item[(c)] \emph{Turn signs.} Consequently the principal turn at each such
corner, a function of that ordered pair, is equal in the half and in $P_2$ ---
in magnitude and in sign.
\item[(d)] \emph{The contact corner.} The single corner $M$ of $L^{*}$ is
replaced by exactly two corners: the contact corner of $L_1$ in $\lambda_1$
and the contact corner of $L_2$ in $\lambda_2$, the ones
Lemma~\ref{lem:contactturnboth} computes. In particular
$c(L_1)+c(L_2)=c(L^{*})+1$.
\item[(e)] \emph{Compatibility.} The bijection of carriers matches each
non-distinguished carrier with a half carrier having the \emph{same} corners
with the same ordered direction pairs and turn signs --- (a)--(c) restricted to
it --- so matched carriers have equal corner counts, equal turn-sign lists, and
equal uniformity status; and $L^{*}$ corresponds to the pair $(L_1,L_2)$ with
(a)--(d).
\end{enumerate}
\end{lemma}
\begin{proof}
$T$ is eligible, so $T\subseteq V_A\sqcup V_B$
(Lemma~\ref{lem:orderedproduct}) and every smoothing site lies strictly inside
$A$ or strictly inside $B$. The newborns are not smoothed, so the traversal
runs uninterrupted through $x_0y_0$ --- the stretch carrying the corner $M$ ---
and through the two remote-edge visits. Write the word as in \eqref{eq:contactword}. Because no newborn is smoothed, the
successor map moves through the four newborn visits without reconnecting: from
$x_0$ it continues to $y_0$ and from $y_0$ into $A$; along the remote stretch it
passes the two remote-edge newborn visits and continues into $B$; and from the
end of $B$ it returns to $x_0$. So the arcs $x_0y_0$ and the remote stretch lie
on one and the same successor cycle, together with whichever arcs of $A$ and of
$B$ the smoothing of $T_A$ and $T_B$ leaves attached to them. \emph{Interval subclaim.} Let an open interval of a cyclic traversal carry finitely many chords,
each with both visits inside it and with no two visit pairs interlacing. Smoothing all of them
leaves exactly one open successor strand joining the two boundary half-edges; every other
successor cycle is internal, with all its arcs strictly inside the interval.

Induct on the number of chords. With none, the interval itself is the one boundary strand and
there is no internal cycle. For the step, order the chords by first visit and choose an
\emph{innermost} one, with no other chord having both visits between its two; one exists because
finiteness and noncrossing leave only nesting. Smoothing it first replaces the sub-arc between
its visits by a closed internal cycle --- the sub-arc contains no other chord visit --- and leaves
the ambient strand entering and leaving where it did, with the two visits spliced out. The
remaining chords lie on that strand or in internal cycles. Those on the ambient strand form a
smaller noncrossing family on an interval with the same boundary half-edges, so induction
applies; a chord in an internal cycle splits it into two internal cycles and does not touch the
boundary strand. Thus exactly one boundary strand remains and every new cycle is internal.
Nested chords give the innermost-first recursion; the empty interval is the base case.
Reconnections at distinct chord pairs act on disjoint visit half-edges and therefore commute,
so simultaneous smoothing may be evaluated in this innermost-first order.

Apply the subclaim first to $A$. The chords of $T_A$ lie wholly inside $A$ and are pairwise
noncrossing there: eligibility makes the support independent, and interlacement of two chords
whose four visits lie in one interval is exactly their crossing there. Hence smoothing $T_A$
leaves one open strand joining the two boundary half-edges of $A$ and otherwise only internal
cycles. The same argument for $B$ and $T_B$ gives its unique boundary strand and internal
cycles. These two boundary strands are exactly the arcs attached to the $x_0y_0$ stretch and
the remote stretch, so they belong to the distinguished cycle, while every remaining cycle is
internal to its interval. Hence the cycles are exactly those inside $A$, those inside $B$, and
the single cycle carrying both stretches, namely $L^{*}$. An earlier revision asserted the
internal-cycle claim in one sentence; the induction above carries it.

Cutting the polygon at the contact replaces the two stretches by the two
closing sub-segments of Definition~\ref{def:halves}: the arcs of $L^{*}$ lying
in $A$ close up into a carrier $L_1$ of $\lambda_1$, and those in $B$ into a
carrier $L_2$ of $\lambda_2$, because the cut separates exactly the two
stretches and the smoothing of $T_A$, $T_B$ is unchanged. The other cycles are
untouched and become carriers of the half they lie in.

For the five exported clauses. (a): every non-contact corner is a vertex or a
smoothing site strictly inside one interval, the newborns being unsmoothed and
$M$ the only corner on the contact stretches; the cut leaves the interior of
each interval pointwise fixed, so the corner lies on the one half carrier that
its cycle becomes. (b): both arcs incident to such a corner lie strictly
inside its interval --- a corner is not a boundary half-edge --- and the cut
alters the curve only at the two contact stretches, so the incoming and
outgoing directions are unchanged as an ordered pair. (c): the principal turn
is a function of that ordered pair, so it is unchanged in magnitude and sign.
(d): the corner $M$ of $L^{*}$ lies on the cut itself; the two closing
sub-segments of Definition~\ref{def:halves} give $L_1$ and $L_2$ one new
corner each at the contact, the ones Lemma~\ref{lem:contactturnboth} computes,
and no other corner is created, the closing sub-segments being straight. Hence
$c(L_1)+c(L_2)=(c(L^{*})-1)+2$. (e): a non-distinguished carrier is untouched
by the cut --- its cycle lies inside one interval --- so its matched half
carrier has literally the same corner set with the same pairs and turns, by
(a)--(c); equal turn-sign lists give equal uniformity status
(Definition~\ref{def:wind}); and $L^{*}$'s data pass to $(L_1,L_2)$ by
(a)--(d).
\end{proof}

\begin{lemma}[minimum carrier corners and selector identities]\label{lem:selectorid}
\begin{enumerate}
\item[(A)] Under the guards of Definition~\ref{def:generic}(A), every carrier of every
support has at least three corners.
\item[(B)] Here eligibility is as in Definition~\ref{def:eligible}.
Let $T$ be eligible, with parts $T_A,T_B$. Write $P_{\mathrm L}$ and
$P_{\mathrm H}$ for the low and high chambers of the event --- upright
subscripts, since the italic $P_H$ already names the polynomial of a residual
piece $H$ (Corollary~\ref{cor:groupedknot}(B)), and an earlier revision wrote
$\wind_{P_H}$ with the two collided. Write $W_T=\wind_{P_{\mathrm H}}(T)$,
$W_1=\wind_{\lambda_1}(T_A)$ and $W_2=\wind_{\lambda_2}(T_B)$, each with its
argument and its ambient polygon. Let $s_{\mathrm{low}}$ be the sign \eqref{eq:S6} evaluated with $P_-$ taken to
be the \emph{low} chamber of the event --- the sign of
$\det(p_b-p_a,\,p_M-p_a)$ there --- which is a datum of the event and of the
chamber, bound here and not borrowed from a lemma printed nine hundred lines
later. Then
\[
\epsilon=1\ \Longrightarrow\ W_T=-s_{\mathrm{low}}\,W_1W_2,
\qquad
\epsilon=0\ \Longrightarrow\ W_T\,W_1W_2=0 .
\]
Moreover, at $\epsilon=1$, the contact carrier $L^{*}$ is uniform if and only if
both half contact carriers $L_1,L_2$ are, and then all three are uniform with
the same turn sign. That clause is proved below with the identities and is
stated here because a consumer reads it as a statement --- an earlier revision
left it inside this proof and cited the lemma for it.

The sign here is $s_{\mathrm{low}}$, fixed by the chamber, and not the
\eqref{eq:S6} sign $s$, which depends on which chamber the parametrisation calls
$P_-$; an earlier revision wrote a bare $s$.
\end{enumerate}
\end{lemma}
\begin{proof}
\emph{Clause (A).} A carrier is a closed curve made of straight segments joined at corners of
nonzero turn. With two corners it would consist of two straight segments
joining the same two points, so the two segments would overlap, which puts an
endpoint of one on the other and vanishes a $\mathrm{G2}$. With one corner it
would contain a straight return of zero length, which is not a segment of a
generic polygon. With no corner at all it would be a closed curve made of a
single straight segment, whose two endpoints coincide, so that segment has zero
length --- excluded for the same reason, and by
Definition~\ref{def:regular}(B)'s requirement that every edge be nonzero. An
earlier revision omitted this case. So there are at least three. In particular a zero selector
means mixed turns, and no one- or two-corner convention is needed.

\emph{Clause (B).} By Lemma~\ref{lem:contactsplit} the carriers other than $L^{*},L_1,L_2$ match
one for one and lie in one half each, with the same corners and hence the same
weights. Write $Q$ for the product of those weights, so that
\[
   W_T=Q\cdot\mathrm{wt}(L^{*}),
   \qquad
   W_1W_2=Q\cdot\mathrm{wt}(L_1)\mathrm{wt}(L_2).
\]
$Q$ is \emph{not} cancelled --- it may be zero, and an earlier revision
cancelled it --- so what is proved below is the local identity between
$\mathrm{wt}(L^{*})$ and $\mathrm{wt}(L_1)\mathrm{wt}(L_2)$, and both displayed
identities follow by multiplying it by $Q$, which is legitimate whatever $Q$ is.
It remains to compare $\mathrm{wt}(L^{*})$ with
$\mathrm{wt}(L_1)\mathrm{wt}(L_2)$. Each of $L_1,L_2$ has at least three
corners (clause~(A), now stated above rather than quoted from
inside another lemma's proof), hence at least two besides its
contact corner.

\emph{Case $\epsilon=1$.} By Lemma~\ref{lem:contactturnboth} all three contact
turns --- the full one and the two half ones --- have the sign of
$s_{\mathrm{low}}$. The non-contact corners of $L^{*}$ partition between $L_1$
and $L_2$ with their signs. So $L^{*}$ is uniform if and only if all its
non-contact corners have the sign of $s_{\mathrm{low}}$, if and only if $L_1$
and $L_2$ are both uniform, and then all
three are uniform of the same sign. If they are mixed, both sides of the
local identity are $0$. If the common sign is right, then $s_{\mathrm{low}}=-1$
and all three weights are $+1$, so
$\mathrm{wt}(L^{*})=1=-s_{\mathrm{low}}\mathrm{wt}(L_1)\mathrm{wt}(L_2)$. If it
is left, then $s_{\mathrm{low}}=+1$, $\mathrm{wt}(L^{*})=(-1)^{c(L^{*})}$ and
$\mathrm{wt}(L_1)\mathrm{wt}(L_2)=(-1)^{c(L_1)+c(L_2)}=(-1)^{c(L^{*})+1}$, so
$\mathrm{wt}(L^{*})=-\mathrm{wt}(L_1)\mathrm{wt}(L_2)
=-s_{\mathrm{low}}\mathrm{wt}(L_1)\mathrm{wt}(L_2)$. Multiplying by $Q$ gives
$W_T=-s_{\mathrm{low}}W_1W_2$.

\emph{Case $\epsilon=0$.} Now the two half contact turns have the sign of
$s_{\mathrm{low}}$ and the full contact turn the sign of $-s_{\mathrm{low}}$.
Suppose $\mathrm{wt}(L^{*})\neq0$; then $L^{*}$ is uniform, necessarily of the
sign of its contact turn, i.e.\ of $-s_{\mathrm{low}}$, so all its non-contact
corners have that sign. Each half carrier inherits at least one such corner and
also carries its own contact corner, of sign $s_{\mathrm{low}}$; so each is
mixed and $\mathrm{wt}(L_1)\mathrm{wt}(L_2)=0$. Symmetrically, if that product
is nonzero then both halves are uniform of sign $s_{\mathrm{low}}$, so every
non-contact corner has that sign, and $L^{*}$ carries those together with its
contact corner of sign $-s_{\mathrm{low}}$, hence is mixed and
$\mathrm{wt}(L^{*})=0$. Either way the local product vanishes, and multiplying
by $Q^2$ gives $W_TW_1W_2=0$.
\end{proof}

\begin{remark}[what the noninterlacing identity replaces]\label{rem:noscalarfact}
Lemma~\ref{lem:selectorid}, clause~(B) does \emph{not} say that in the noninterlacing
branch the full selector factors through the half selectors with a sign. It
says the opposite: a live full selector and a live product of half selectors
are mutually exclusive. Any argument that replaces the full selector by the
half-selector product in that branch is refuted by that clause, not supported by
it.
\end{remark}

\subsection{The interlacing endpoint}\label{sec:sr12}

\begin{lemma}[what the two chambers of a bigon event share]\label{lem:chambercommon}
Here eligibility is as in Definition~\ref{def:eligible} and $\rot$ is as in
Definition~\ref{def:rot}.
Let $T$ be eligible at a simple bigon event. Write $C_{\mathrm L}$ for its
contact carrier in the low chamber and $C_{\mathrm H}$ for the same carrier in
the two-crossing chamber. Then, with no hypothesis on any selector:
\begin{enumerate}
\item[(i)] every carrier of $T$ other than the contact carrier is the same
closed curve, with the same corners, the same turn signs and the same residual
pieces, in the two chambers; write $U_T$ for the product of the reads
$\Omega_1$ over those carriers, one number for both chambers;
\item[(ii)] $C_{\mathrm L}$ and $C_{\mathrm H}$ have the same corners and the
same turn signs, hence $\wind(T)$ takes one value $W_T$ in both chambers, and
$\rot(C_{\mathrm L})=\rot(C_{\mathrm H})$, hence
$R(C_{\mathrm L})=R(C_{\mathrm H})$.
\end{enumerate}
\end{lemma}
\begin{proof}
(i) The two chambers differ by moving the single vertex $M$ across the remote
edge, and the newborn crossings lie in the bigon disc that move sweeps. A
carrier not through the contact contains neither $M$ nor either newborn, so
neither its corner list nor its crossing list changes; the reads over them
therefore agree, and their product is one number $U_T$.

(ii) The two copies of the contact carrier differ by a Reidemeister-II finger.
Neither newborn is a polygon vertex and neither lies in $T$, so neither is a
corner; the corner list is therefore the same list of polygon vertices and
smoothing sites, and each corner's turn is a determinant of two edge directions
whose sign is guarded by (G1) or (G5) and constant on the event's interval
(Lemma~\ref{lem:guardconst}). Hence the weights, and $\wind(T)$, agree. During a standard finger
deformation this fixed cyclic corner list has nonzero edges and no consecutive
negative-multiple pair: at every corner its two directions have the guarded
nonzero determinant just identified. It is therefore a path in the corresponding
$\mathcal R_c$, so Lemma~\ref{lem:turnlift}(ii) keeps $\rot$ unchanged.
\end{proof}

\begin{definition}[the marked data of an eligible support]\label{def:markeddata}
Under Convention~\ref{conv:s7standing}, let $T$ be eligible with parts
$T_A,T_B$, and let
$L_1,L_2$ be the contact carriers of $T_A$ in $\lambda_1$ and of $T_B$ in
$\lambda_2$ (Lemma~\ref{lem:contactsplit}). Write
\[
   f_i(a)=\bigl[z^0\bigr]Q_{i}(a,z),\qquad
   d_i=1-w_i-R_i,\qquad
   \omega_i=\bigl[a^{d_i}\bigr]f_i,\qquad
   D=d_1+d_2,
\]
where $Q_i$ and $w_i$ are the grouped polynomial and grouped writhe of $L_i$,
$\rot_i=\rot(L_i)$ its signed rotation and $R_i=|\rot_i|$, so that $\Omega_1$ on
$L_i$ is $\omega_i$. The signed rotation carries one name in this document,
$\rot$ with a subscript --- numerals for halves, upright $\mathrm L,\mathrm H$
for chambers; earlier revisions also wrote $\rho_i$ and $r_i$ for the same
quantity, which is three names for one thing. The chambers and the half selectors are bound here as
well, since both are used outside the lemma that first introduced them:
$P_{\mathrm L}$ and $P_{\mathrm H}$ are the low and the high chamber of the
event --- upright subscripts, the italic $P_H$ being the polynomial of a
residual piece $H$ --- and
\[
   W_1=\wind_{\lambda_1}(T_A),\qquad W_2=\wind_{\lambda_2}(T_B)
\]
are the two half selectors, each with its argument and its ambient half. An
earlier revision left them bound only inside Lemma~\ref{lem:selectorid}, clause~(B) and used
them in several statements after it. The letter $Q$ is used for a half's grouped polynomial and
$F_\nu$ for a chamber's, $\nu\in\{\mathrm L,\mathrm H,\mathrm A\}$. For the
symbols bound in this definition --- $f_i$, $d_i$, $\omega_i$, $Q_i$, $w_i$,
$\rot_i$, $R_i$, $W_i$, $L_i$ --- a numerical subscript means a half and an
upright $\mathrm L$ or $\mathrm H$ means a chamber. The rule is stated for those
symbols and not as a general convention, because chambers elsewhere do carry
numerical names: $P_0$ and $P_2$ are the no-crossing and two-crossing chambers
of a vertex-on-edge event, and $d_0,d_1$ are the incident wall tangents of the
sliding branch. An earlier revision wrote "numerical subscripts mean halves and
never chambers", which those three names contradict.

\emph{The chamber data.} For $\nu\in\{\mathrm L,\mathrm H\}$ write $C_\nu$ for
the contact carrier of $T$ in that chamber and $D_\nu$ for the grouped diagram
that carrier bears (Corollary~\ref{cor:groupedknot}(B)), and put
\[
   w_\nu=w(D_\nu),\quad R_\nu=R(C_\nu),\quad \rot_\nu=\rot(C_\nu),\quad
   d_\nu=1-w_\nu-R_\nu,\quad F_\nu=P_{D_\nu},\quad
   \Omega_\nu=\bigl[a^{d_\nu}z^0\bigr]F_\nu,
\]
the writhe of that diagram, the absolute and signed rotations of its curve, the
slot, the diagram's polynomial and its read. The quantities are defined
\emph{from the diagrams} $D_{\mathrm L},D_{\mathrm H}$, which is what the
full-twist lemma quantifies over; an earlier revision defined them from the
chambers and left the diagrams to be inferred. These six symbols are used throughout the sections
below and were bound only in running prose, which is not a binder.

Two factors are bound \emph{separately}, an earlier revision having named only
their product:
\[
   W_T=\wind_{P_{\mathrm H}}(T)=\wind_{P_{\mathrm L}}(T)
   \quad\text{(the two agree by Lemma~\ref{lem:chambercommon}(ii))},
\]
and $U_T$ the product of the reads $\Omega_1$ over the carriers of $T$ other
than the contact carrier --- the untouched ones, common to the two chambers by
Lemma~\ref{lem:chambercommon}(i), which is the clause about those carriers;
clause (ii) is about the contact carrier and the selector, and an earlier
revision cited it here. The term of $T$ is then $W_TU_T$ times the read on the contact
carrier.
\end{definition}

\begin{remark}[why the hypothesis is the selector alone]\label{rem:sr12hyp}
The hypothesis of Lemma~\ref{lem:triplebridge}(vii) is $W_T\neq0$ and not $W_TU_T\neq0$. The proof below uses only
that the full carrier through the contact is live, which is what $W_T\neq0$
gives; the spectator product $U_T$ appears nowhere in it. The recorded defect was actual: an earlier revision's statement asked for
$W_TU_T\neq0$, while its consumer had only $W_T\neq0$; the invocation therefore
exceeded the stated hypothesis --- the peer's point. The repair weakens that
statement to $W_T\neq0$, which its proof consumes, and the live full-twist
calculation below consumes exactly that repaired interface.

\end{remark}

\subsection{The sliding branch}\label{sec:sliding}

Now let the event be of sliding type: the two neighbours $p=p_{M-1}$ and
$q=p_{M+1}$ lie strictly on \emph{opposite} sides of the line of the remote
edge. Convention~\ref{conv:s7standing} is in force throughout this subsection: the
restart changes the branch, not the standing hypotheses, and (H3) is what makes
the halves of this branch polygons that $X_1$ is defined on. Write $r=p_b-p_a$,
$d_0=p_M-p$, $d_1=q-p_M$ for the three directed wall tangents at the event.

\begin{lemma}[signed-word transport, the containing-sector Cartesian product, and the chamber-to-half carrier map]
\label{lem:slidingcartesian}
\begin{enumerate}
\item[(A)] For this statement, the \emph{signed cyclic Gauss word} of a generic polygon
means its Gauss word (Section~\ref{sec:setup}) together with, at each double
point, the record of which visit is the overpass under the divide convention of
Definition~\ref{def:piecediagram}. Every divide crossing is positive, so this
is the whole decoration; there is no separate crossing-sign decoration.

Let $x_-$ be the crossing present in $P_-$ and $x_+$ the one present in $P_+$,
and let $\phi$ relabel $x_-$ as $x_+$ and fix every other label. Then $\phi$ is
an isomorphism of those signed cyclic Gauss words
$W(P_-)\to W(P_+)$. Consequently it
is an isomorphism of interlacement graphs, a bijection of supports, and it
matches smoothing cycles, undominated sets, residual components, their signed
piece diagrams, their polynomials, their writhes, and the carrier owning each
component.
\item[(B)] Let the event be a simple sliding vertex-on-edge event, let
$P_\bullet\in\{P_-,P_+\}$ be either chamber, let $x_\bullet$ be the relocated
crossing in it --- $x_-$ in $P_-$ and its image $x_+$ in $P_+$ --- and write
$G=G_{P_\bullet}$ and $G_{\lambda_1},G_{\lambda_2}$ for the interlacement graphs
of that chamber and of the two halves at the wall. Nothing below distinguishes
the two chambers: the statement is proved once for either, and an earlier
revision stated it for $P_-$ alone while its consumers quantified over both.
Cutting the traversal circle of $P_\bullet$ at the two visits of $x_\bullet$
separates the old labels into two intervals;
the induced subgraph of $G$ on the labels of one interval \emph{is} the
interlacement graph of the corresponding half, since interlacement of two chords
with all four visits inside one interval is read from their order inside it,
which the half inherits. With that identification,
\[
   \{\,S\in\Ind(G):x_\bullet\in S\,\}
   \;\longleftrightarrow\;
   \Ind(G_{\lambda_1})\times\Ind(G_{\lambda_2}),
   \qquad
   S\longmapsto\bigl(S_1,S_2\bigr),
\]
is a bijection, where $S_i=(S\setminus\{x_\bullet\})\cap V(G_{\lambda_i})$.
\item[(C)] Fix a simple sliding event, a chamber $P_\bullet\in\{P_-,P_+\}$, and a support
$S\in\Ind(G_{P_\bullet})$ containing the relocated crossing. Let $S_1,S_2$ be its
halves (Lemma~\ref{lem:slidingcartesian}(B)). Let $\chi_\bullet$ send a carrier $C$
of $S$ in $P_\bullet$ to the successor cycle of the half $\lambda_i$ carrying
the same retained visits, $i$ being the interval $C$ meets. Then $\chi_\bullet$
is a well-defined map from the carriers of $S$ in $P_\bullet$ to the carriers of
$S_1$ in $\lambda_1$ and of $S_2$ in $\lambda_2$, and it has the following
properties.
\begin{enumerate}
\item[(i)] \emph{Correspondence.} Every carrier of $S$ in $P_\bullet$ meets
exactly one of the two intervals, and $\chi_\bullet$ is a bijection onto the
union of the two halves' carrier sets. Exactly two carriers pass through the
contact, one in each interval, and they are sent to the two halves' contact
carriers, one each. Only one of the two carries the short collapsing edge, and
it is the one in the interval containing that edge.
\item[(ii)] \emph{Records.} For each carrier $C$, the retained-visit map from
$C$'s marked points to those of $\chi_\bullet(C)$ is an orientation-preserving
record isomorphism (Definition~\ref{def:record}).
\item[(iii)] \emph{Owners and grouped data.} The residual pieces of $S$ carried
by $C$ are exactly the residual pieces of the corresponding half support carried
by $\chi_\bullet(C)$; each has the same label set, hence the same piece diagram,
polynomial and writhe. Thus the matched carriers have the same grouped polynomial
and grouped writhe.
\item[(iv)] \emph{Corners.} Away from the contact, $C$ and $\chi_\bullet(C)$
have the same corner list with the same signs. Of the two carriers through the
contact, the one \emph{not} carrying the short collapsing edge also has the same
corner list as its image; the one carrying it has, at the contact, one direction
more in its local direction list than its image has, and nothing is claimed here
about the two lists beyond that --- the refinement is computed where it is
consumed.
\item[(v)] \emph{Compatibility with the chamber map.} If $S$ contains the
relocated crossing of $P_-$ and $\phi$ is the chamber-to-chamber relabelling,
then $\chi_+\circ\phi=\chi_-$ as maps from the carriers of $S$ in $P_-$ to the
halves' carriers.
\end{enumerate}
\end{enumerate}
\end{lemma}
\begin{proof}
\emph{Clause (A).}
\emph{Source.} This proof is a transcription of the peer's pack \texttt{evidence/codex/phase15-sliding-universal-attack/REPORT.md}, SHA-256 \texttt{4f6da162\dots9318}, §``Signed word and support transport'',
lines 118--150. The argument is the peer's and not this author's; the display
numbering $\mathrm{U}$ is that document's.

The two visits of the relocated crossing occupy two local intervals of the
cyclic traversal word, and event isolation keeps every old visit out of those
intervals. Sliding replaces the visit on one incident edge by the visit on the
other, so deleting the relocated label leaves the same old cyclic word.

The opposite-side condition supplies the signed refinement: both directed
incident tangents point into the same open half-plane bounded by the directed
remote tangent, which is
\begin{equation}
   \sgn\det(r,d_0)=\sgn\det(r,d_1)\neq0 .
   \tag{U8}
\end{equation}
So the determinants against the remote tangent have the same sign, and under
the positive-crossing convention the over/under role and the oriented sign of
the gained visit are those of the lost visit. Hence $\phi$ is an isomorphism of
the \emph{signed} cyclic word and not only of its unsigned chord diagram; it
induces an isomorphism of interlacement graphs and a bijection of all
independent supports.

The smoothing successor permutation is determined by the signed traversal word
together with the selected labels, so corresponding supports have corresponding
successor cycles. The undominated residual graph is preserved with them, hence
so are its connected components. For the piece data the map is restricted
rather than waved at: for a residual piece $H$, $\phi$ carries the visits of $H$
bijectively to those of $\phi H$, preserving their cyclic order --- $\phi$ being
an isomorphism of the cyclic word --- and preserving the over/under designation
and the sign at each, by the positive-divide clause above. Those are the four
conditions of Definition~\ref{def:record}, so $\phi|_H$ is a record isomorphism
between the two piece diagrams' records, and Axiom~\ref{ax:gausscode}
identifies the links: $P_H=P_{\phi H}$ and $|H|=|\phi H|$. A residual component
is owned by the successor cycle containing its visits, so carrier ownership is
preserved as well. Every combinatorial and polynomial assertion of the statement follows. Rotation
is not among them: this lemma claims nothing about it, and the carrier rotations
are the subject of a separate lemma below.

\emph{Clause (B).}
\emph{Source.} A transcription of the peer's pack \texttt{evidence/codex/phase15-sliding-universal-attack/REPORT.md}, SHA-256 \texttt{4f6da162\dots9318}, §``Carrier rotation'', lines 165--170, where it is
Display~U6.

A support chord is independent from the relocated chord $x_\bullet$ exactly when
both of its ends lie in one half interval, since a chord with one end in each
interval interlaces $x_\bullet$. Two chords whose ends lie in \emph{different}
half intervals cannot interlace each other, their visits being in disjoint arcs.
Hence, \emph{for a set of old labels each of which is nonadjacent to
$x_\bullet$}, independence in the old graph is independence of its two halves in
the respective half graphs.

That hypothesis is not decoration and the statement is false without it: in the
word $123123$ the chords $1,2,3$ pairwise interlace, so a set may fail to be
independent while its two halves --- one of them empty --- are independent in
their own graphs. The bijection of this lemma is unaffected, because the sets it
maps are supports \emph{containing} $x_\bullet$, and independence of such a
support already forces every other member to be nonadjacent to $x_\bullet$,
which is the hypothesis. An earlier revision stated the equivalence for every
set of old labels. Conversely any pair of independent supports, one from
each half, combines with the relocated crossing to an independent support
containing it. The two assignments are mutually inverse.

\emph{Clause (C).}
Cut the traversal circle of $P_\bullet$ at the two visits of the relocated
crossing $x_\bullet$. By Lemma~\ref{lem:slidingcartesian}(B) the old labels split
into the two half intervals and the induced subgraphs are the half interlacement
graphs, so $S_i$ is a support of $\lambda_i$.

\emph{The two intervals close separately.} $S$ contains $x_\bullet$, so
$x_\bullet$ \emph{is} smoothed. An oriented smoothing at a self-crossing of one
closed curve joins each outgoing branch to the incoming branch that continues
the orientation, and so cuts the traversal circle into exactly two cycles, one
carrying the labels of each interval. Every other element of $S$ is independent
from $x_\bullet$ and therefore has both visits inside one interval
(Lemma~\ref{lem:slidingcartesian}(B)), so no later smoothing joins the two
families. Hence every carrier meets exactly one interval, and smoothing $S$
inside an interval performs exactly the reconnections that smoothing $S_i$
performs in $\lambda_i$; the successor cycles correspond one for one, which is
the bijection of (i). Two of them run through the contact, one per interval,
and the short collapsing edge lies in exactly one interval, hence on exactly one
of those two.

An earlier revision said instead that the contact carrier was ``the one cycle
that meets both intervals''. That is the picture of a \emph{bigon} event, where
the newborns are not smoothed and one cycle does meet both; here the relocated
crossing is smoothed and it is never true. The minimal support $S=\{x_\bullet\}$
is the sharpest case: nothing else is smoothed, and the two intervals close into
two carriers, each mapping to the contact carrier of its half. Those carriers
need not be pieceless --- a minimal support forbids further \emph{smoothings},
not residual pieces, and the half may carry pieces of its own --- and an earlier
revision called them pieceless. The
selector identity below already reads the two-carrier picture --- it multiplies
one factor $c$ for the chamber's short-edge carrier by one factor $h$ for the
other --- so the correction removes an inconsistency rather than creating one.

For (ii): the correspondence is the identity on the retained visits, it
preserves their cyclic order by Lemma~\ref{lem:carrierword} applied inside each
interval, and it preserves the over/under designation and the sign at each,
those being read from the two branch directions, which the cut does not move.
For (iii) the label sets must first be shown to agree, and an earlier revision
assumed it here and proved it later from this very clause. The argument is
three lines and belongs here. Remove the relocated crossing from $S$ and
intersect what remains with the two half label sets, which is the definition of
$S_1,S_2$; every label interlacing the relocated crossing is dominated and so
absent from $U(S)$; every remaining label has both visits in one interval, and
there is no interlacement edge between the two intervals, a chord with ends in
different intervals interlacing the relocated chord and so not being present. So
\[
   G_{P_\bullet}\bigl[U(S)\bigr]=G_{\lambda_1}\bigl[U(S_1)\bigr]\sqcup
                     G_{\lambda_2}\bigl[U(S_2)\bigr],
\]
and the residual components on the two sides are the same sets of labels. Then
Lemma~\ref{lem:carriers}(iv), applied in $\lambda_i$, puts each component on
one carrier, and Definition~\ref{def:piecediagram} gives it the same piece
diagram, the label set and its signs being unchanged; the passage from the equal
records to equal polynomials is Axiom~\ref{ax:gausscode}, which is what licenses
reading a polynomial off a record at all. For (iv): a corner of a carrier is a polygon vertex or a smoothing site
(Definition~\ref{def:wind}), and away from the contact each is present in
$C$ and in $\chi_\bullet(C)$ with the same incoming and outgoing directions,
since the wall moves only the sliding vertex; so the corner lists and their
signs agree there. The carrier through the contact that does not carry the short
edge meets the contact in a single corner, present in both with the same
directions, so its whole list agrees too. The remaining carrier traverses the
short edge, which the wall contracts to a point: its local direction list at the
contact has the two directions of the collapsing corner where the image has
one, which is the one extra direction claimed, and no more is asserted here.

For (v): $\phi$ is the identity on the old labels and carries $x_-$ to $x_+$, so
the two chambers' traversal circles have the same cyclic word on the labels of
$S$ and $\phi S$, the same cut, and the same reconnections. The successor cycles
therefore correspond under $\phi$, and each pair is sent to the same half
carrier, since the half data are read from the same interval word in both. That
is $\chi_+\circ\phi=\chi_-$.
\end{proof}

\begin{remark}[what this map replaces]\label{rem:chamberhalfwhy}
Three arguments below --- rotation transport, coefficient-product transport and
selector transport --- previously each reached for the halves in its own way: by
``restricting'' the chamber-to-chamber relabelling $\phi$, which is a map
between the two \emph{chambers} and not between a chamber and a half, or by
citing Proposition~\ref{prop:sectortarget}(A), which is stated for the both-newborn support
of a \emph{bigon} event and not for a sliding one. The missing object is one and the
same map, and it is printed here.
\end{remark}

\begin{lemma}[rotation transport, sliding incidence, the collapsing half, coefficient products, selector transport and the sliding jump]\label{lem:slidingcoeff}
\begin{enumerate}
\item[(A)] Here $\rot$ is as in Definition~\ref{def:rot}.
Under $\phi$, corresponding carriers have the same signed rotation and the same
$R$. More precisely:
\begin{enumerate}
\item[(i)] for a support \emph{avoiding} the relocated crossing, every carrier
of one chamber is a regular homotopy of its mate, and they have the same signed
rotation, the same $R$, \emph{and the same corner-sign list};
\item[(ii)] for a support \emph{containing} it, every carrier $C$ of either
chamber has the signed rotation and $R$ of its wall-half mate
$\chi_\bullet(C)$; hence $\chi_+\circ\phi=\chi_-$ gives the two chambers the same
signed and absolute rotations, carrier by carrier. No claim is made about
corner-sign lists in this sector: the carrier carrying the short collapsing edge
has one corner more than its wall-half mate, so the lists have different lengths.
\end{enumerate}
\item[(B)] Under Convention~\ref{conv:s7standing}, let the event be of sliding type, write
$r=p_b-p_a$ for the direction of the remote edge and $d_0=p_M-p_{M-1}$,
$d_1=p_{M+1}-p_M$ for the two incident edge directions at the moving vertex. These are the three directed wall tangents, restated
here rather than imported by citation.
\begin{enumerate}
\item[(i)] At a simple sliding event, $\sgn\det(r,d_0)=\sgn\det(r,d_1)\neq0$. Moreover
there is an interval of the event path around the wall on which both statements
below hold, and they are about different objects. First, each incident
\emph{line} meets the \emph{line} of the remote edge at a point whose parameter
along the remote \emph{segment} lies strictly inside $(0,1)$ --- that is a
condition on the remote segment, and it holds for both incident lines. Second,
exactly one incident \emph{edge} meets the remote edge in the relative interiors
of both, in each of the two chambers, and the two chambers exchange which
incident edge that is; the other incident edge has its line meeting the remote
segment at a point beyond its own far endpoint. No other crossing changes on the
interval. An earlier revision compressed the two into one sentence that read
"each incident line meets the remote line at a parameter strictly inside $(0,1)$
along both segments", which asks the intersection to lie inside the incident
segment as well and so contradicts the "exactly one" it was supporting.

The interval restriction is not a formality. Opposite sides of the remote
\emph{line} do not by themselves put the intersection inside the remote
\emph{segment}: the segments $(0,0)$--$(1,0)$ and $(2,1)$--$(2,-1)$ have
endpoints strictly on opposite sides of the first one's line and meet that line
at $(2,0)$, outside the first segment. The peer's example is that one. Since the
incidence parameter is continuous and equals the contact point's at the wall,
which lies strictly inside the remote segment by
Convention~\ref{conv:events}, shrinking the interval puts it inside for every
parameter of the interval. The statement is phrased symmetrically in the two
chambers rather than assigning a direction, which the event's parametrisation
does not fix.
\item[(ii)] Let $s$ be the sign
\eqref{eq:S6} taken in the initial chamber. Put
\[
   \eta_1=\sgn\det(r,d_1),\qquad \eta_2=\sgn\det(d_0,r).
\]
Then $\eta_2=-\eta_1$; in each chamber exactly one half carries the short
collapsing edge; and that half has contact-corner sign $s$ in the initial
chamber and $-s$ in the other.

The last clause is \emph{true} and is proved below from the smoothing successor.
What an earlier revision got wrong was the \emph{reason}, not the conclusion: it
inferred the owner from which side of the remote line the incident edge lies on,
and the peer's chart refutes that inference --- with the remote edge
$(-\tfrac12,0)\to(\tfrac12,0)$, $M=(0,0)$, $p=(0,-1)$, $q=(1,1)$ and
$M_\mp=(0,\pm\tfrac1{10})$, the half with $\eta_1=s$ is $q$'s while $q$ lies on
the \emph{same} side as $M_-$. A later revision then over-corrected and printed
that the conclusion itself is false, which it is not. Ownership is read from the
smoothing successor. Smoothing the relocated crossing splits the traversal at its two
visits into two cyclic words, and the short edge --- the stretch of the incident
edge between the moving vertex and the contact --- lies in exactly one of them,
namely the one whose arc contains the moving vertex; that arc is the half whose
contact corner is the turn at the contact on that side. There are exactly two
cases, according to which incident edge carries the crossing in the chamber
under consideration, and clause~(B)(i) says the two chambers
exchange them. The caveat recorded there applies here too: opposite sides of the
remote \emph{line} do not place an intersection inside the remote
\emph{segment}, so the event interval is shrunk as clause~(B)(i) prescribes before
any of this is read off.
\end{enumerate}
\item[(C)] Let $S\in\Ind(G_{P_-})$ be arbitrary, and let $\phi S$ be its image under the
relabelling of Lemma~\ref{lem:slidingcartesian}(A). Write $A_-(S)$ for the product
of the coefficient factors $\Omega_1$ over the carriers of $S$ in the initial
chamber and $A_+(\phi S)$ for the same product over the carriers of $\phi S$ in
the other chamber.
\begin{enumerate}
\item[(i)] One always has $A_-(S)=A_+(\phi S)$. If $x_-\in S$, put
$S_i=(S\setminus\{x_-\})\cap V(G_{\lambda_i})$ and write $A_i(S_i)$ for the
product of the coefficient factors over the carriers of $S_i$ in $\lambda_i$.
Every argument is bound: an earlier revision wrote $A_i$ with none. Under that condition,
Lemmas~\ref{lem:slidingcartesian}(C)(i)--(iii) and~\ref{lem:slidingcoeff}(A)(ii) identify the two
chamber products carrier by carrier with the two half products, and
\[
   A_-(S)=A_+(\phi S)=A_1(S_1)A_2(S_2).
\]
\item[(ii)] If $x_-\in S$, then, under Convention~\ref{conv:s7standing}, put
$S_i=(S\setminus\{x_-\})\cap V(G_{\lambda_i})$, the two halves of
Lemma~\ref{lem:slidingcartesian}(B), and write
$W_-=\wind(S)$, $W_+=\wind(\phi S)$ and $W_i=\wind(S_i)$, each with its
argument. With $r,d_0,d_1$ the three directed wall tangents as above, put
\[
 \eta_1=\sgn\det(r,d_1),\qquad \eta_2=\sgn\det(d_0,r),\qquad
 \tau=\sgn\det(d_0,d_1).
\]
Here $\eta_2=-\eta_1$, and the collapsing half has contact-corner sign $s$ in
the initial chamber and $-s$ in the other, where $s$ is the sign~\eqref{eq:S6}
in the initial chamber. For $\eta\in\{\pm1\}$, write $h_\eta$ for the weight of
the distinguished carrier of the half with contact-corner sign $\eta$, and
$c_\eta$ for the weight of the corresponding chamber carrier carrying the short
edge, whose contact-corner sign there is $\eta$. Write $Q$ for the product of
the weights of all carriers not distinguished at the contact. It is the same
number in both chambers and both half terms: Lemma~\ref{lem:slidingcartesian}(C)(i)
matches those carriers bijectively, and clauses (iv) and (v) give the matched
corner-sign lists. An earlier revision wrote $W_-,W_+,W_1,W_2$ and $Q$ with no
arguments and no bijection. Then
\[
   W_1W_2=Q\,h_s h_{-s},\qquad
   W_-=Q\,c_s h_{-s},\qquad
   W_+=Q\,h_s c_{-s},\qquad
   c_\eta=\begin{cases}-\eta\,h_\eta,&\tau=\eta,\\ 0,&\tau=-\eta,\end{cases}
\]
and consequently $W_+-W_-=s\,W_1W_2$.
\end{enumerate}

\item[(D)] At every simple sliding vertex-on-edge event, in the sense of
Convention~\ref{conv:events}, whose two halves are generic polygons,
\[
   X_1(P_+)-X_1(P_-)=s\,X_1(\lambda_1)X_1(\lambda_2),
\]
with $s$ the sign \eqref{eq:S6} in the initial chamber and $\lambda_1,\lambda_2$
the halves of Definition~\ref{def:halves}.
\end{enumerate}
\end{lemma}
\begin{proof}
\emph{Clause (A).}
\emph{Source.} This proof is a transcription of the peer's pack \texttt{evidence/codex/phase15-sliding-universal-attack/REPORT.md}, SHA-256 \texttt{4f6da162\dots9318}, §``Carrier rotation'', lines 152--218.

\emph{(i), the avoiding sector.} Every carrier corner is either a polygon
vertex or the smoothing site of a persistent transverse crossing. Its point and
its incoming and outgoing directions vary continuously through the wall, and all
corner determinants remain nonzero by the guards: the vertex turn by (G1), the
smoothing-site turn by the active (G5) of the crossing, which persists. The
contact itself is only a transverse coincidence of parameter values, the
relocated crossing not being smoothed here. Thus each successor cycle has a fixed cyclic corner list. The nonzero
corner determinants make every incident edge nonzero and exclude consecutive
negative multiples, so it is a path in the corresponding $\mathcal R_c$.
Lemma~\ref{lem:turnlift}(ii) keeps its signed rotation constant; taking
magnitudes gives $R$. The same nonzero
determinants preserve the entire corner-sign list, which is also the selector
statement in this sector.

\emph{Three corners.} Every smoothed carrier has at least three corners under
the generic straight-edge guards: a carrier with two corners would consist of
two straight segments joining the same two points, which would overlap, and one
with a single corner would contain a zero-length straight return. This matters
because a zero selector then means mixed turns rather than a special one- or
two-corner convention.

\emph{(ii), the containing sector.} The correspondence with the halves is
Lemma~\ref{lem:slidingcartesian}(C)(i), and clause~(iv) of that lemma says the corner
lists agree away from the contact. There are \emph{two} carriers through the
contact, one in each interval, and clause (iv) gives the one not carrying the
short collapsing edge the same corner list as its mate as well; so every carrier
except the single one carrying that edge is a regular homotopy of its mate, by
the argument of (i). An earlier revision spoke of ``the one through the
contact'', which is the picture Lemma~\ref{lem:slidingcartesian}(C)(i) withdraws. The
distinguished carrier has one short edge that contracts
at the wall, and its local direction list is one of the two refinements
\begin{equation}
   (r,d_0,d_1)\longrightarrow(r,d_1),
   \qquad
   (d_0,d_1,r)\longrightarrow(d_0,r),
   \tag{U9}
\end{equation}
the two-direction lists on the right being the corresponding wall-half corners.
Let $\theta$ denote the principal oriented turn between two nonparallel
directions. Display~(U8) places both incident directions in one open angular
interval of length $\pi$ bounded by the remote direction, so no branch cut is
crossed and the two local angle additions
\begin{equation}
   \theta(r,d_0)+\theta(d_0,d_1)=\theta(r,d_1),
   \qquad
   \theta(d_0,d_1)+\theta(d_1,r)=\theta(d_0,r)
   \tag{U10}
\end{equation}
are exact, with no hidden full turn. All other turns agree by the polygon
path of (i). Lemma~\ref{lem:turnlift}(ii) therefore gives every chamber
carrier the same signed rotation and,
after taking absolute values, the same $R$ as its wall-half mate; clause~(v)
then makes the two chambers agree in both values carrier by carrier.

The pack also gives an equivalent form that avoids the angle bookkeeping, and it
is worth printing because it is the one a formalizer would prefer: choose one
reference ray avoiding the finite set of wall carrier directions and shrink the
chamber interval until that same ray avoids all corresponding directions, then
compute every rotation with that common ray. On either refinement in~(U9) a ray
in a backtracked subarc is crossed once with each sign and cancels, and every
other ray gets the direct-corner count.

\emph{Clause (B).} (i) Both incident tangents are transverse to $r$: a parallel one would put its far
endpoint on the remote line, an incidence outside the zero set
Convention~\ref{conv:events} allows, which contains no incidence but the
contact one. The neighbours lie on opposite sides of
the line, so the traversal $p\to M\to q$ crosses the line at $M$ and the two
incident directions point into the same open half-plane bounded by $r$, i.e.\
the two determinants have one sign. The statement is invariant: for each of the two incident edges, that edge meets
the remote edge exactly when its two endpoints lie strictly on opposite sides of
the remote line, and $p$, $q$ lie on opposite sides of it while $M$ passes from
one side to the other. So in each chamber exactly one incident edge meets the
remote edge, and the two chambers exchange which one it is --- no direction of
travel is named. No other crossing changes: any other change would vanish a
member of $\mathcal G$ outside the zero set, and by
Lemma~\ref{lem:guardconst} every relevant member outside $Z$ is nonzero of
constant sign on a possibly smaller interval, on which the whole argument is
read.

(ii) By clause~(B)(i) the two incident directions satisfy
$\sgn\det(r,d_0)=\sgn\det(r,d_1)$, and $\det(d_0,r)=-\det(r,d_0)$, so
$\eta_2=-\eta_1$: the two wall-half contact turns are opposite.

For the assignment the reason is the \emph{cross-pairing} of the oriented
smoothing, and not a side of a line. An earlier revision of this proof concluded
by the far-side inference that the statement above already withdraws --- the
statement was corrected and the proof under it was not, which is a way of
leaving a refuted argument in print.

At a slide exactly one incident edge meets the remote segment
(clause~(B)(i)), so in the chamber under consideration the
relocated crossing sits on one of them; write $X$ for the crossing point. Take
first the incoming edge $pM$. The traversal visits $X$ twice, once running
$p\to M$ and once running along the remote edge $a\to b$. The oriented smoothing
joins the branch arriving along $pM$ to the branch leaving along $X\to b$, and
the branch arriving along $a\to X$ to the branch leaving along $X\to M$. So the
short stretch $X\to M$ continues into $M\to q$: it lies on the successor cycle
carrying $q$'s arc, which is the half $\lambda_1$, whose contact corner is
$\eta_1$.

Take instead the outgoing edge $Mq$, which is the case in the other chamber by
clause~(B)(i). Now the two visits of the crossing $Y$ run $M\to q$
and $a\to b$, the smoothing joins the branch arriving along $M\to Y$ to
$Y\to b$ and the branch arriving along $a\to Y$ to $Y\to q$, and the short
stretch $M\to Y$ is entered from $p$: it lies on the cycle carrying $p$'s arc,
the half $\lambda_2$, whose contact corner is $\eta_2=-\eta_1$.

Both orientations occur and both are printed, an earlier revision of this proof
having treated only the first. In the initial chamber the crossing sits on the
incoming edge exactly when the moving vertex has passed to the side where $pM$
meets the remote segment; then the short stretch lands on $\lambda_1$, whose
contact corner is $\eta_1$, and $\eta_1=s$ there. In the reversed configuration
--- the peer's chart, with $M_-=(0,-\tfrac1{10})$ and $s=-1$, where $M_-q$
crosses the remote edge at $(\tfrac1{11},0)$ --- the crossing sits on the
outgoing edge instead, the short stretch lands on $\lambda_2$, and $\eta_2=s$.
The two cases are exchanged by the chamber, by clause~(B)(i), so
in either configuration the collapsing half has contact-corner sign $s$ in the
initial chamber and $-s$ in the other, which is the assignment the consumers
read.

What carries the argument in both cases is which branch the smoothing attaches
the short stretch to. That is a statement about the successor, and not about
which side of the remote line anything lies on.

\emph{Clause (C).} (i) Lemma~\ref{lem:slidingcartesian}(A) preserves the grouped polynomials and the
grouped writhes carrier by carrier, and Lemma~\ref{lem:slidingcoeff}(A) preserves
the absolute rotations, so each factor
$\Omega_1=[a^{1-w-R}z^0]P$ is selected at the same index from the same
polynomial: the first identity follows factor by factor. When $x_-\in S$, the
Lemmas~\ref{lem:slidingcartesian}(C)(i)--(iii) and~\ref{lem:slidingcoeff}(A)(ii) identify
the carriers with the disjoint union of the two sets of half carriers and
preserve each owner, grouped polynomial, grouped writhe and $R$, so the product
over the mapped carriers splits into the two half products. No coefficient
extraction is distributed over a polynomial product: each carrier lies in one
half.

(ii) By clause~(B)(ii) the collapsing half is the one of
contact-corner sign $s$ in the initial chamber and the one of sign $-s$ in the
other; that is what makes the three displayed products the $\pm$ pair they are.
In the chamber in which a given half carries the collapsing short edge, its
single wall corner of sign $\eta$ is refined into two corners, of signs $\eta$
and $\tau$. If the wall carrier is mixed, both $h_\eta$ and $c_\eta$ are $0$,
since adding a corner cannot remove an existing opposite turn. If it turns
right only and $\tau$ agrees, the refined carrier still turns right only and
its weight is $1=-\eta h_\eta$ with $\eta=-1$ and $h_\eta=1$. If it turns left
only and $\tau$ agrees, the refined carrier has one more corner, so its weight
is negated: $c_\eta=-h_\eta=-\eta h_\eta$ with $\eta=+1$. If $\tau$ disagrees
the refined carrier is mixed and $c_\eta=0$. Those are all cases, proving the
displayed rule. The three product formulas are the definition of $\wind$ with
the distinguished factor named. The single $Q$ is legitimate because
Lemma~\ref{lem:slidingcartesian}(C)(i), (iv) and (v) identifies the other carriers and
their corner-sign lists in both chambers and both half terms. Finally, if
$\tau=s$ then $c_s=-s\,h_s$ and $c_{-s}=0$, so
$W_+-W_-=Q(0+s\,h_sh_{-s})=s\,W_1W_2$; if $\tau=-s$ then $c_s=0$ and
$c_{-s}=s\,h_{-s}$, so $W_+-W_-=Q(s\,h_sh_{-s}-0)=s\,W_1W_2$.

\emph{Clause (D).} Split the supports of each chamber into those avoiding the relocated crossing
and those containing it. On the avoiding sector the coefficient products agree by
clause~(C)(i), and the selectors agree by
Lemma~\ref{lem:slidingcoeff}(A)(i), whose corner-sign list is exactly the selector
datum: no carrier acquires or loses a corner there. So that sector cancels term
by term. Clause~(C)(ii) is \emph{not} invoked here; it is
stated for a support containing the relocated crossing and an earlier revision
of this proof cited it in the avoiding sector, outside its own quantifier. On the containing sector, write the difference as the sum over the supports
themselves before any reindexing:
\[
   \sum_{\substack{S\in\Ind(G_{P_-})\\ x_-\in S}}
   \bigl(W_+(\phi S)-W_-(S)\bigr)\,A_-(S),
\]
which is legitimate because $\phi$ is a bijection of the containing supports of
the two chambers (Lemma~\ref{lem:slidingcartesian}(A)) and $A_-(S)=A_+(\phi S)$
(clause~(C)(i)), so the two chambers' containing sectors are
summed over one index set. An earlier revision passed straight from the jump to
the Cartesian sum, omitting this line, which is the first equality of the
source's final display. Now
clause~(C)(i) makes the non-selector part of the term the product
$A_1(S_1)A_2(S_2)$ of the two half terms' non-selector parts. The reindexing of
the sum by $(S_1,S_2)\in\Ind(G_{\lambda_1})\times\Ind(G_{\lambda_2})$ is
Lemma~\ref{lem:slidingcartesian}(B); the carrier, owner and piece-data
identification at a fixed $S$ is Lemma~\ref{lem:slidingcartesian}(C)(i)--(iii), and
Lemma~\ref{lem:slidingcoeff}(A)(ii) supplies its rotation data. Finally
clause~(C)(ii) makes the selector difference $s\,W_1W_2$.
Hence
\[
X_1(P_+)-X_1(P_-)
=\sum_{S_1,S_2}s\,W_1(S_1)A_1(S_1)\,W_2(S_2)A_2(S_2)
=s\,X_1(\lambda_1)X_1(\lambda_2),
\]
where $A_i$ denotes the coefficient product of the half term.
\end{proof}

\begin{remark}[printed from the source, and what it replaces]\label{rem:slidingsource}
This is the rotation argument of the counterpart's sliding pack, printed in its
own terms: regular homotopy for the avoiding sector, and exact angle addition
under the opposite-side condition for the contracting edge. An earlier revision
of this document argued it instead by an axiom about embedded arcs in a disc,
which the peer broke --- for the empty support the carrier meets a contact disc
in two arcs, not one, and they cross. That axiom is not used here, and the
displayed additions are what make the absence of a hidden full turn a
computation rather than an assurance.
\end{remark}

\subsection{Universal (S7)}\label{sec:s7universal}


\begin{remark}[the two sectors of the noninterlacing case]\label{rem:nistatus}
Four structural points of Section~\ref{sec:nijoin} are recorded here so that
no reader has to reconstruct them. The orientation factor:
Lemma~\ref{lem:nitransport}, clause~(A) names $s$
and says that the $s=-1$ configuration is the reflection of the $s=+1$ one,
with every determinant negating together, instead of printing $+1$ globally and
consuming $s$ later. The $U_T=0$ branch: the assembly multiplies by $U_T$
whatever its value, so a vanishing spectator product vanishes the term. The
same-piece hypothesis of the transport lemma: it is the source's own and it is
a hypothesis of that lemma. The
\emph{different-piece sector}: its common-neighbourhood input is
Lemma~\ref{lem:contactword}(ii); its discharge is printed in
Lemma~\ref{lem:chambertransport}, clauses~(A) and~(B),
Theorem~\ref{thm:s7universal}(D)(i), Theorem~\ref{thm:s7universal}(D)(iii) and
Theorem~\ref{thm:s7universal}(D)(v); Theorem~\ref{thm:s7universal}, clause~(E) does not exclude it.

\emph{The two sectors, named.} For an eligible $T$ at a bigon event, the
undominated set $U(T)$ carries the induced interlacement graph
$G_{P_2}[U(T)]$, and the two newborns $x,y$ lie in it. The \emph{same-piece}
sector is the set of eligible $T$ for which $x$ and $y$ lie in one connected
component of that graph, and the \emph{different-piece} sector is the set for
which they lie in two; the names are the residual pieces those components
become, and every eligible $T$ is in exactly one sector. An earlier revision of
this remark used the two names without defining them.

\emph{The different-piece discharge.} All four of its branches are proved outright, and
the branch statements are read at $\epsilon=0$ --- the noninterlacing branch,
which is where the different-piece sector arises, the interlacement bit of
Definition~\ref{def:sectors} being $0$ exactly when the two newborn chords do
not interlace. With $\epsilon=0$ fixed: the two newborns are singleton positive
kinks (Lemma~\ref{lem:chambertransport}, clause~(A), whose hypothesis that is); the row
$f_\nu=[z^0]F_\nu$ of the contact carrier in chamber $\nu$, and its prefactor,
transport between the chambers with no selector hypothesis
(Lemma~\ref{lem:chambercommon} and Lemma~\ref{lem:chambertransport}, clause~(B)); the
one-newborn terms vanish by cutting the affected carrier at its surviving
singleton (Theorem~\ref{thm:s7universal}(D)(i) and Theorem~\ref{thm:s7universal}(D)(iii)); and the
high-neither and low terms read the degrees $D-4$ and $D-2$, in the writhes
$w_{\mathrm L},w_{\mathrm H}$ of Definition~\ref{def:markeddata}, which the
floor on the two halves kills --- the same two degrees, and the same floor, as
Theorem~\ref{thm:s7universal}, clause~(C) uses in the same-piece sector
(Theorem~\ref{thm:s7universal}(D)(v)).
\end{remark}

\subsection{The returned aggregate in bracket form}\label{sec:bracket}

This subsection derives, for the interlacing branch, the bracket-form
identification of the returned aggregate. It closes two of its three terms;
what remains after it is one displayed identity, stated at the end.

\begin{lemma}[rotation across the contact, both branches]\label{lem:rotshift}
Here eligibility is as in Definition~\ref{def:eligible} and $\rot$ is as in
Definition~\ref{def:rot}.
Let $L^{*}$ be the contact carrier of an eligible $T$, let $L_1,L_2$ be its two
halves, and let
\[
   s_{\mathrm{low}}=\sgn\det\bigl(p_b-p_a,\;p_M-p_a\bigr)\ \text{ in the
   \emph{low} chamber }P_0
\]
be the side of the remote line on which the moving vertex sits before the
crossing is born. Then
\[
   \rot(L_1)+\rot(L_2)=\rot(L^{*})+\begin{cases}0,&\epsilon=1,\\
   s_{\mathrm{low}},&\epsilon=0.\end{cases}
\]
The sign here is $s_{\mathrm{low}}$ and not the sign $s$ of \eqref{eq:S6},
which is read in whichever chamber the parametrisation calls $P_-$. The two are
related by the orientation of the path: $s=\delta\,s_{\mathrm{low}}$, where
$\delta=+1$ if the parametrisation calls the low chamber $P_-$ and $\delta=-1$
if it calls the high chamber $P_-$ --- a definition, not a claim, since $s$ and
$s_{\mathrm{low}}$ are the same determinant read in the two chambers and that
determinant changes sign across the wall. The statement above is invariant
under
reparametrisation because $s_{\mathrm{low}}$ is defined by the chamber and not
by the direction of travel.
\end{lemma}
\begin{proof}
The three carriers have the same non-contact corners with the same turns
(Lemma~\ref{lem:contactsplit}), so by Lemma~\ref{lem:turnlift}(ii)
\[
   2\pi\bigl(\rot(L_1)+\rot(L_2)-\rot(L^{*})\bigr)=\eta_1+\eta_2-\tau,
\]
where $\eta_1,\eta_2$ are the two half contact turns and $\tau$ the full one.
The chart of Lemma~\ref{lem:contactturnboth} is fixed by placing the moving
vertex on a chosen side of the remote line, which is the choice
$s_{\mathrm{low}}=+1$; in that chart $\eta_1=\arg q$ and $\eta_2=\pi-\arg p$, while $\tau$ is the principal
value of $\arg q-\arg p+\pi$. At $\epsilon=1$ that value lies in $(0,\pi)$ and
equals $\arg q-\arg p+\pi$, so $\eta_1+\eta_2-\tau=0$. At $\epsilon=0$ it lies
in $(-\pi,0)$ and equals $\arg q-\arg p-\pi$, so $\eta_1+\eta_2-\tau=2\pi$ and
the shift is $+1$.

The configuration with $s_{\mathrm{low}}=-1$ is the image of the one with
$s_{\mathrm{low}}=+1$ under a reflection of the plane, which reverses every signed rotation and every turning
angle. So at $\epsilon=1$ the shift $0$ is unchanged, while at $\epsilon=0$ the
shift $+1$ becomes $-1$. Writing the two cases as the single symbol
$s_{\mathrm{low}}$ gives the display.
\end{proof}

\begin{remark}[the shift carries the side sign]\label{rem:rotshiftsign}
An earlier revision printed the constant $+1$ at $\epsilon=0$. That is the
$s=+1$ half of the statement, and the peer refuted the constant form on the
exact guarded $24$-gon of Remark~\ref{rem:diffcensus}, whose low-chamber side
sign is $s_{\mathrm{low}}=-1$
and whose signed rotations are $\rot(L_1)=\rot(L_2)=1$ and $\rot(L^{*})=3$, so
that $\rot(L_1)+\rot(L_2)-\rot(L^{*})=-1$. A second revision bound the shift to the
\eqref{eq:S6} sign $s$, which is read in whichever chamber is called $P_-$ and
therefore flips with the parametrisation while the geometry does not; the
binder is now the chamber-defined
$s_{\mathrm{low}}$. Lemma~\ref{lem:nitransport}, clause~(B) consumes the corrected form; no
conclusion drawn from it changes.
\end{remark}

\begin{remark}[this is the missing full turn]\label{rem:missingturn}
Lemma~\ref{lem:rotshift} is the reason naive addition of principal angles
across the contact is wrong in the noninterlacing branch, which the phase-1.5
record warns of without displaying the discrepancy: there the two halves carry
one full turn more than the carrier they came from, and at $\epsilon=1$ they
carry exactly as much.
\end{remark}

\begin{remark}[a withdrawn route]\label{rem:withdrawnbracket}
An earlier revision of this document printed here a lemma asserting that after
polarization the contact carrier is, in the low chamber, the connected sum of
the two half groupings, with framed row $\Gamma_L=\Gamma_1\Gamma_2$; a
proposition deriving the high-neither minus low difference from it by the
framed self-skein; and a corollary reducing the interlacing branch to an
aggregate over selector-dead supports. \textbf{All three are withdrawn.} The first
is false, and its falsity is visible inside this document: the writhe ledger
$w(D_-)=w_1+w_2+2\ell-1$ of Lemma~\ref{lem:triplebridge}(v) at $\epsilon=1$,
for a diagram whose
two parts are joined by $\ell$ between-interval crossings, so
$\Gamma_L=\Gamma_1\Gamma_2$ would force $2\ell=1$. The error was to think that
polarizing the between-interval crossings into a descending order lets
Reidemeister~II separate the two parts. It does not: in a divide every crossing
is positive, so between-interval crossings all have one sign and form a twist
region, which no Reidemeister~II move removes. The correct identity carries the
framing factor, $\Gamma_L=a^{2\ell-1}\Gamma_1\Gamma_2$, and the surviving
source factors the two parts only \emph{after} a contact crossing is smoothed,
retaining the single carrier's framing datum before that.

What this costs is the whole of my route to the interlacing endpoint. What
remains from it is stated where it belongs: Lemma~\ref{lem:rotshift}, the
rotation shift, which is independent of the framing question; and
Theorem~\ref{thm:s7universal}, clause~(B)(i) below, which is a determinant argument in
the contact chart and uses none of the withdrawn material.
\end{remark}

\begin{remark}[two lemmas retired, and one debt discharged]
\label{rem:excretired}
This subsection printed, in earlier revisions, two lemmas about the one-newborn
branches at $\epsilon=1$ --- that at most one of them is live, and that a
uniform contact carrier kills both --- and a remark calling the surviving
branch's coefficient read ``the remaining debt of the whole bigon branch''. All
three are withdrawn, and none is replaced by a weaker claim.

The peer refuted the two lemmas as printed: both treated liveness as one sign on
every raw corner of $A$ and one sign on every raw corner of $B$, whereas after a
smoothing the corners belong to \emph{carriers} and uniformity is carrierwise.
The statement that makes them unnecessary is the peer's.
Theorem~\ref{thm:s7universal}, clause~(B)(i) kills \emph{both} one-newborn branches for
\emph{every} eligible support, uniform contact carrier or not, by an exact
determinant argument in the printed chart. With both branches always dead,
``at most one is live'' is vacuous, and the debt is discharged rather than
outstanding: there is no surviving branch whose coefficient read anyone must
compute.

The consumers are rewired to Theorem~\ref{thm:s7universal}, clause~(B)(i).
\end{remark}

\subsection{The full-twist calculation}\label{sec:fulltwist}

This subsection prints the interlacing join as one live-support calculation.
The statement below carries its own $\epsilon=1$ and eligible-support binder and
imports all chamber and half data from Definition~\ref{def:markeddata}. Its
proof chooses $q=x$ and binds $D_{\mathrm A},F_{\mathrm A}$ before invoking
Lemma~\ref{lem:triplebridge}(ii); it binds $\ell$ only after that clause supplies
the ordered two components. No rotation, target or $\Omega$ is attached to the
two-component auxiliary link. Lemma~\ref{lem:fulltwist} remains branch-neutral.

\begin{lemma}[the abstract full-twist triple]\label{lem:fulltwist}
Let $D_{\mathrm L},D_{\mathrm H},D_{\mathrm A}$ be oriented diagrams and let $q$
be a positive crossing of $D_{\mathrm H}$ such that
\begin{enumerate}
\item[(T1)] the oriented smoothing of $D_{\mathrm H}$ at $q$ is
$D_{\mathrm A}$;
\item[(T2)] switching $q$ in $D_{\mathrm H}$ gives a diagram carried to
$D_{\mathrm L}$ by oriented Reidemeister-II moves.
\end{enumerate}
The symbols this lemma reads off a diagram are bound here, from the diagrams it
quantifies over, and are not imported from any definition attached to a polygon
or a chamber. They are bound in two groups, because they have two domains. For
\emph{any} oriented diagram $D$ --- any number of components --- put
\[
   F_D=P_D ,
\]
the HOMFLY--PT polynomial of the link it presents (Axiom~\ref{ax:homfly}). For
a diagram $D$ carried by a \emph{single} closed plane curve $\Gamma_D$, and
only for such a diagram, put
\[
   d(D)=1-w(D)-R(\Gamma_D),\qquad
   \Omega(D)=\bigl[a^{d(D)}z^0\bigr]F_D ,
\]
with $w(D)$ the writhe and $R(\Gamma_D)=|\rot(\Gamma_D)|$ the absolute
rotation of that curve. An earlier revision bound all three on one domain,
which typed $d$ and $\Omega$ at a two-component diagram, where the rotation
number of Definition~\ref{def:rot} is not defined. Abbreviate $F_\nu=F_{D_\nu}$ for $\nu\in\{\mathrm L,\mathrm H,\mathrm A\}$,
and $d_\nu=d(D_\nu)$, $\Omega_\nu=\Omega(D_\nu)$ for
$\nu\in\{\mathrm L,\mathrm H\}$ \emph{only}. The restriction is not tidiness:
$D_{\mathrm A}$ is the oriented smoothing of $D_{\mathrm H}$ and is carried by a
\emph{two-component} link in the application, while $R(\Gamma_D)$ is the
absolute rotation of one closed curve (Definition~\ref{def:rot}), so $d$ and
$\Omega$ are not defined at $\mathrm A$ at all. Nothing below asks for them:
$D_{\mathrm A}$ enters only through $F_{\mathrm A}$ and through the coefficient
$[a^{d_{\mathrm L}-1}z^{-1}]F_{\mathrm A}$, which is a coefficient of that
polynomial and not a read at a slot of its own. An
earlier revision used $d_\nu$ and $\Omega_\nu$ in this statement while binding
them only in Definition~\ref{def:markeddata}, which is about the contact
carriers of an eligible support; the abstract triple has no support, and the
lemma was not readable on its own hypotheses. The two bindings agree wherever
both apply, that definition computing $d_\nu$ and $\Omega_\nu$ from these same
grouped diagrams. Then
\[
   F_{\mathrm H}=a^{-2}F_{\mathrm L}+a^{-1}zF_{\mathrm A}.
\]
If moreover $d_{\mathrm H}=d_{\mathrm L}-2$, then
\[
   \Omega_{\mathrm H}-\Omega_{\mathrm L}
   =\bigl[a^{d_{\mathrm L}-1}z^{-1}\bigr]F_{\mathrm A}.
\]
\end{lemma}
\begin{proof}
Apply the skein of Axiom~\ref{ax:homfly} at $q$: its $L_+$ is $D_{\mathrm H}$,
its $L_-$ is the switched diagram, whose polynomial is $F_{\mathrm L}$ by (T2)
and the Reidemeister-II invariance of the polynomial, and its $L_0$ is
$D_{\mathrm A}$ by (T1). So $aF_{\mathrm H}-a^{-1}F_{\mathrm L}=zF_{\mathrm A}$,
which is the first display.

For the second, $\Omega_{\mathrm H}=[a^{d_{\mathrm H}}z^0]F_{\mathrm H}
=[a^{d_{\mathrm L}-2}z^0]F_{\mathrm H}$. Taking $[a^{d_{\mathrm L}-2}z^0]$ of
the first display, the first term contributes
$[a^{d_{\mathrm L}-2}z^0](a^{-2}F_{\mathrm L})
=[a^{d_{\mathrm L}}z^0]F_{\mathrm L}=\Omega_{\mathrm L}$ and the second
contributes $[a^{d_{\mathrm L}-1}z^{-1}]F_{\mathrm A}$.
\end{proof}

\begin{remark}[why this lemma is branch-neutral]\label{rem:branchneutral}
Nothing in the statement mentions $\epsilon$, a contact, or a polygon: it is an
identity about three oriented diagrams related by one crossing. That is
deliberate. An earlier revision proved the $\epsilon=1$ case and then had the
$\epsilon=0$ section import it ``verbatim'', which is not an argument --- the
two branches have different contact words and different auxiliary components.
Here each branch proves (T1) and (T2) for itself, in
Lemma~\ref{lem:triplebridge}, and the algebra is quoted once. The peer's
blueprint is the source of this arrangement.
\end{remark}

\begin{lemma}[the residual sets split, and who owns them]\label{lem:ownermap}
Here eligibility is as in Definition~\ref{def:eligible}.
Let $T$ be eligible with parts $T_A,T_B$, let $A$ and $B$ be the two intervals
of the contact word of Lemma~\ref{lem:contactword}, let $V_A,V_B$ be the old
labels with both visits in $A$ and in $B$ respectively, and let $\mathcal N$ be
the old labels with one visit in each. Write $U(T)$ for the undominated set of $T$ in the high-neither
chamber and $U_{\lambda_i}(T_i)$ for the undominated set of $T_i$ in the half
$\lambda_i$. Then
\[
   U(T)\cap V_A=U_{\lambda_1}(T_A),\qquad
   U(T)\cap V_B=U_{\lambda_2}(T_B),\qquad
   U(T)\cap\mathcal N=\mathcal N\setminus N_{G_{P_2}}(T),
\]
and every label of the third set has one visit in each interval.

\end{lemma}
\begin{proof}
A chord with both visits in one interval and a chord with both visits in the
other cannot interlace, their visits lying in disjoint arcs
(Lemma~\ref{lem:orderedproduct}); and by Lemma~\ref{lem:contactword}(i) every
old label lies in $V_A$, in $V_B$ or in $\mathcal N$. Let $c\in V_A$. Then $c$
interlaces no element of $T_B$, so $c\in U(T)$ if and only if $c\notin T_A$ and
$c$ interlaces no element of $T_A$; interlacement of two chords with all four
visits inside $A$ is read from their order inside $A$, which the half inherits,
so that is exactly $c\in U_{\lambda_1}(T_A)$. For $c\in V_B$ interchange the
two intervals throughout the previous three sentences: $c$ interlaces no element
of $T_A$, so $c\in U(T)$ if and only if $c\notin T_B$ and $c$ interlaces no
element of $T_B$, which is $c\in U_{\lambda_2}(T_B)$. For $c\in\mathcal N$ the condition $c\in U(T)$ is $c\notin T$ and
$c\notin N_{G_{P_2}}(T)$; the first holds for every eligible $T$ by
Definition~\ref{def:eligible}, so it reduces to $c\notin N_{G_{P_2}}(T)$, which
is the third identity. The subscript is not decoration: $N(T)$ without one is
bound in Definition~\ref{def:branches} as the \emph{scalar} high-neither term of
$T$, and an earlier revision wrote the graph neighbourhood with the same
symbol. The last clause is the definition of $\mathcal N$.

The third identity is the one that must be read carefully, and an earlier
revision did not: a mixed label that is \emph{not selected} may still be
\emph{dominated}, and a dominated label is erased from the diagram rather than
retained. Membership of $U(T)$, not failure of membership of $T$, is what
retention means (Definition~\ref{def:pieces}).
\end{proof}

\begin{remark}[the newborn-free statement against the both-newborn one]\label{rem:ownervssliding}
Lemma~\ref{lem:ownermap} concerns the newborn-free support, not
Proposition~\ref{prop:sectortarget}(B), which concerns $S=T\cup\{x,y\}$ at a noninterlacing bigon. There the support
contains both newborns, every label of $\mathcal N$ is adjacent to them and is
dominated outright, and the residual sets carry no mixed label at all. Here the
support contains no newborn, so a mixed label survives unless an old selection
dominates it, and which ones survive is the content of the third identity.
\end{remark}

\begin{lemma}[different-piece singleton kinks and chamber transport of an eligible row]\label{lem:chambertransport}
\begin{enumerate}
\item[(A)] Here eligibility is as in Definition~\ref{def:eligible}.
Let $\epsilon=0$ and let $T$ be eligible. If the two newborns lie in different
residual pieces of $T$, then those pieces are the singletons $\{x\}$ and
$\{y\}$; each singleton's piece diagram is an unknot, so $P_{\{x\}}=P_{\{y\}}=1$
and $w(\{x\})=w(\{y\})=1$.
\item[(B)] Here eligibility is as in Definition~\ref{def:eligible}.
Under Convention~\ref{conv:s7standing}, let $T$ be eligible at a simple bigon
event, let $C_{\mathrm L}$ and $C_{\mathrm H}$ be its contact carriers in the
low and the high chamber, let $W_T$ be its chamber selector and $U_T$ the
product of the reads over its untouched carriers --- all four as bound in
Definition~\ref{def:markeddata}, which is where they are defined, rather than
imported from the lemma that first computed with them. Then, with no hypothesis on
any selector:
\begin{enumerate}
\item[(iii)] if $T$ is different-piece, the residual pieces of $C_{\mathrm H}$ are those
of $C_{\mathrm L}$ together with the two singletons $\{x\}$ and $\{y\}$, so the grouped
rows agree, $f_{C_{\mathrm H}}=f_{C_{\mathrm L}}$, while the writhes differ by two,
$w_{T,C_{\mathrm H}}=w_{T,C_{\mathrm L}}+2$.
\end{enumerate}
Consequently --- and this conclusion is stated only under the
different-piece hypothesis of subclause~(iii), since its rows are derived from it ---
writing $U_T$ for the product of the reads on the carriers of
Lemma~\ref{lem:chambercommon}(i) --- clause (B) has no subclause (i) since the
split, and an earlier revision still pointed at one --- and
$d_0=1-w_{T,C_{\mathrm L}}-R(C_{\mathrm L})$,
\[
   L(T)=W_TU_T\bigl[a^{d_0}\bigr]f_{C_{\mathrm L}},
   \qquad
   N(T)=W_TU_T\bigl[a^{d_0-2}\bigr]f_{C_{\mathrm L}},
\]
both expressions being valid whether or not $W_TU_T$ vanishes.
\end{enumerate}
\end{lemma}
\begin{proof}
\emph{Clause (A).} At $\epsilon=0$ the newborns do not interlace, so they are not adjacent. Suppose
$x$ had an undominated old neighbour $c$. By Lemma~\ref{lem:contactword}(ii) $c$ is a
neighbour of $y$ as well, so $x,c,y$ is a path in the residual graph
$G_{P_2}[U(T)]$ and the two newborns share a piece, contrary to hypothesis.
Hence $x$ has no residual neighbour at all and $\{x\}$ is a component of the
residual graph. Every hypothesis used --- eligibility of $T$, $\epsilon=0$, and
Lemma~\ref{lem:contactword}(ii), whose conclusion is an equivalence between the two
newborns --- is invariant under interchanging the names $x$ and $y$, and so is
the conclusion that the two lie in different pieces. Interchanging them
therefore yields the identical derivation with $y$ in place of $x$, and
$\{y\}$ is a component as well.

For the diagram: by Definition~\ref{def:piecediagram} the diagram of the piece
$\{x\}$ is the parent curve with \emph{every} double point other than $x$
erased and $x$ resolved by the divide convention. A closed curve with exactly
one crossing is an unknot diagram, so its normalized polynomial is $1$
(Definition~\ref{def:homfly}), and its writhe is $w(\{x\})=|\{x\}|=1$ because
the divide convention makes every crossing positive.

\emph{Clause (B).} (iii) By clause~(A) the two newborns are singleton
residual pieces of $T$ in the high chamber; every other residual piece consists
of old crossings, and its diagram erases $x$ and $y$, so the two strands there
pass without interaction and the finger is undone by a planar isotopy, leaving
the same diagram as in the low chamber. Each singleton contributes the factor
$1$ to the grouped polynomial and the summand $1$ to the grouped writhe
(clause~(A)), which is the displayed statement.

The two displays are then Definition~\ref{def:X1} read in each chamber, the
slot in the high chamber being
$1-w_{T,C_{\mathrm H}}-R(C_{\mathrm H})=1-(w_{T,C_{\mathrm L}}+2)-R(C_{\mathrm L})=d_0-2$. Nothing above divides by a
selector or a spectator read, so both displays hold as identities of integers
whatever those factors are.
\end{proof}

\begin{lemma}[the vertex-on-edge triple, endpoint bracket, auxiliary components and live full-twist jump]
\label{lem:triplebridge}
Here eligibility is as in Definition~\ref{def:eligible}.
Let $T$ be eligible at a simple bigon event with interlacement bit $\epsilon\in\{0,1\}$.
Clauses (i)--(v) also require $\epsilon=1$, or $\epsilon=0$ with $T$ same-piece. This statement
carries its own data and inherits none: let $D_{\mathrm L}$ and $D_{\mathrm H}$
be the grouped diagrams of the contact carrier in the two chambers
(Corollary~\ref{cor:groupedknot}(B)); let $q:=x$ be the newborn crossing whose
first visit opens the contact word of Lemma~\ref{lem:contactword}; let
$D_{\mathrm A}$ be the oriented smoothing of $D_{\mathrm H}$ at $q$; and let
$F_\nu=P_{D_\nu}$ for $\nu\in\{\mathrm L,\mathrm H,\mathrm A\}$. The symbol
$\ell$ is \emph{not} bound here: its binding needs the two components that
clause~(ii) below establishes, and a conclusion cannot retroactively type a
prior let --- an earlier revision bound "the linking number of the two
components, when it has two", a partial let discharged only by its own lemma.
An earlier revision also took $F_{\mathrm A}$ and $\ell$ from a preamble that
fixes $\epsilon=1$, while the statement covers both branches. The choice of newborn is a \emph{binding} and not a
convenience: the assignment of writhes to components in (iii) is not equivariant
under exchanging the two newborns, and an earlier revision quantified $q$ over
both while proving the statement for this one. Then
\begin{enumerate}
\item[(i)] $q$ is a positive crossing of $D_{\mathrm H}$; the oriented smoothing
of $D_{\mathrm H}$ at $q$ is $D_{\mathrm A}$; and switching $q$ in
$D_{\mathrm H}$ gives a diagram carried to $D_{\mathrm L}$ by oriented
Reidemeister-II moves. These are the three hypotheses that
Lemma~\ref{lem:fulltwist} asks of an abstract triple, restated here rather than
named, so that this clause can be read without opening that lemma;
\item[(ii)] $D_{\mathrm A}$ has exactly two components, ordered by the
half-label convention --- the first is the one containing the first half's
labels, the second the other; their grouped polynomials are $Q_1$ and $Q_2$,
the two half polynomials of Definition~\ref{def:markeddata}.

\emph{With the two components of (ii) so ordered, let $\ell$ be their linking
number; clauses (iii)--(v) read it.}

\item[(iii)] the second component has writhe $w_2$, and the first has writhe
$w_1+(1-\epsilon)$;
\item[(iv)] $[z^{-1}]F_{\mathrm A}=a^{-2\ell}(a-a^{-1})f_1f_2$;
\item[(v)] $w_{\mathrm L}=w_1+w_2+2\ell-\epsilon$.
\item[(vi)] without the extra hypothesis imposed on clauses (i)--(v),
\[
   w_{\mathrm H}=w_{\mathrm L}+2,\qquad R_{\mathrm H}=R_{\mathrm L},\qquad
   d_{\mathrm H}=d_{\mathrm L}-2.
\]

\item[(vii)] With the marked-data symbols as in Definition~\ref{def:markeddata},
if $\epsilon=1$ and $W_T\neq0$, then
\[
   \bigl[a^{D-2}\bigr]f_1f_2-\bigl[a^{D}\bigr]f_1f_2+\omega_1\omega_2=0 .
\]
A summand with $W_T=0$ contributes zero to the weighted aggregate whatever its
bracket, so the weighted SR.12 aggregate vanishes.

\item[(viii)] Under Convention~\ref{conv:s7standing}, if $\epsilon=1$, bind the parts $T_A,T_B$ and the following chamber and
half data exactly as in Definition~\ref{def:markeddata}:
$C_\nu,D_\nu,F_\nu,w_\nu,\rot_\nu,R_\nu,d_\nu,\Omega_\nu$ for
$\nu\in\{\mathrm L,\mathrm H\}$; $L_i,Q_i,f_i,w_i,\rot_i,R_i,d_i,\omega_i$
for $i\in\{1,2\}$; and $D,W_T$. If $W_T\neq0$, then
\[
   \Omega_{\mathrm H}-\Omega_{\mathrm L}=-\omega_1\omega_2.
\]
\end{enumerate}
\end{lemma}
\begin{proof}
\emph{(i).} (T1) holds by the definition of $D_{\mathrm A}$. For (T2), switch
$q$: the two newborns then carry opposite signs and, the contact disc being
empty (Convention~\ref{conv:events}), they bound an empty bigon face, so an
oriented Reidemeister-II move deletes both and leaves the low grouped diagram
$D_{\mathrm L}$ --- the old pieces and their crossings being untouched by the
move, which is supported in that disc.

\emph{(ii) and (iii).} Read the contact word of Lemma~\ref{lem:contactword}. At
$\epsilon=1$ it is $x_0y_0Ax_1y_1B$; smoothing $q=x$ cuts the traversal at
$x_0,x_1$ and reconnects, so the two visits of $y$ land on \emph{different}
components. Then $y$ is a \emph{mixed} crossing of the two-component link
$D_{\mathrm A}$, and the disposition of a mixed crossing is not that of a
self-crossing: it is not deleted, and Axiom~\ref{ax:homfly} is not what applies
to it. An earlier revision wrote that it was deleted with no change of
polynomial by that axiom, which types it as a self-crossing of one component.
What is true, and what is used below, is that a mixed crossing belongs to
neither component's own diagram, so it enters neither $f_1$ nor $f_2$; it
belongs to the linking data, and enters through the exponent $-2\ell$ that
Lemma~\ref{lem:homflyrows}(ii) consumes in~(iv). At $\epsilon=0$ the word is
$x_0y_0Ay_1x_1B$; smoothing $x$ puts both visits of $y$ on one component, where
it is a positive self Reidemeister-I curl --- there the crossing \emph{is} a
self-crossing, its deletion is a Reidemeister-I move, and the polynomial is
unchanged. In both branches the two components carry, respectively, the
labels of $A$ together with the half $\lambda_1$'s and those of $B$ together
with $\lambda_2$'s.

That the two components' grouped polynomials are $Q_1,Q_2$ is the record
argument and not a picture. The residual-owner identities are
Lemma~\ref{lem:ownermap}'s --- Lemma~\ref{lem:orderedproduct} is only a support
bijection and Lemma~\ref{lem:contactsplit} only a carrier split, and neither
supplies the owner map, which an earlier sentence here attributed to them.
The identities must be read at the level of \emph{labels}, not of pieces. It is \emph{false} that each
old residual piece of $T$ on the contact carrier lies wholly in $A$ or wholly in
$B$, and an earlier revision asserted it; the peer refuted it on the label set
\texttt{contact-0001}, where the piece
$\{(0,2),(0,3),(1,3),(2,7),(3,6)\}$ has one crossing among the first half's
labels, none among the second's, and four whose two visits fall in different
intervals. A piece is a connected component of an interlacement graph and
nothing makes it respect the cut.

What the record isomorphism needs, and what those two lemmas give, is a
statement about crossings. Call a crossing of $D_{\mathrm H}$ \emph{mixed} when
its two visits lie in different intervals of the contact word; by
Lemma~\ref{lem:contactword}(ii) the mixed old labels are exactly $\mathcal N$.
Of those, the ones \emph{retained} in the diagram are
$\mathcal N\setminus N_{G_{P_2}}(T)$
and not all of $\mathcal N$: a mixed label dominated by a selected old crossing
is erased --- dominated meaning adjacent in $G_{P_2}$ to a selected label ---
and Lemma~\ref{lem:ownermap} is what says which survive. An
earlier revision wrote that an eligible $T$ contains no mixed label and
concluded that every mixed crossing is retained, which confuses not-selected
with undominated.

The retained mixed crossings belong to neither component's own diagram of
$D_{\mathrm A}$, each having one visit in each interval: they are the crossings
whose signed count is $2\ell$. Delete them from the two component records and
every remaining crossing has both visits in one interval, hence lies on one
component; by Lemma~\ref{lem:ownermap} the remaining labels of the
component through $A$ are exactly $U_{\lambda_1}(T_A)$, the residual labels of
the first half, and those of the other component exactly
$U_{\lambda_2}(T_B)$. That is the owner map, and it is a statement about
residual \emph{sets}: Lemma~\ref{lem:orderedproduct} is a bijection of supports
and Lemma~\ref{lem:contactsplit} a splitting of carriers, and neither of them
says this. Restricting the traversal
of the smoothed diagram to the retained visits of one component gives a
bijection onto the visits of the corresponding half carrier, and by
Lemma~\ref{lem:carrierword} it preserves their cyclic order; the over/under
designation and the sign at each are those of the parent, by
Definition~\ref{def:piecediagram}, since a smoothing elsewhere alters neither.
So the retained-visit map is an orientation-preserving record isomorphism in the
sense of Definition~\ref{def:record}, and Axiom~\ref{ax:gausscode}, invoked
through Corollary~\ref{cor:groupedknot}(B) on each component, identifies the two
grouped diagrams with the half grouped knots. Their polynomials are $Q_1,Q_2$.
A half carrier bearing no piece is outside the $k\geq1$ scope of
Corollary~\ref{cor:groupedknot}(B) and takes one line of its own: the component has
no retained crossing at all, so its diagram is an embedded circle, its link type
is the unknot, and its grouped polynomial is $1$ with writhe $0$, which is
Theorem~\ref{thm:carrierfloor}(D) and is the value $Q_i$ takes there.

The writhes follow by counting positive crossings: the second component carries
exactly the half $\lambda_2$'s residual crossings, so $w_2$; the first carries
$\lambda_1$'s together with the surviving newborn when $\epsilon=0$, and none
when $\epsilon=1$, so $w_1+(1-\epsilon)$.

\emph{(iv).} $D_{\mathrm A}$ is a two-component link, and it is \emph{not}
claimed to be a split union --- its two components link, and the factor
$a^{-2\ell}$ records that. Lemma~\ref{lem:homflyrows}(ii) applied to it, with the two
component rows $f_1,f_2$ from (ii), gives
$[z^{-1}]F_{\mathrm A}=a^{-2\ell}(a-a^{-1})f_1f_2$.

\emph{(v).} Count writhes, branch by branch, and count them in
$D_{\mathrm H}$, whose crossings are the residual crossings of the two halves,
the mixed ones, and the two newborns $x$ and $y$; all are positive
(Definition~\ref{def:piecediagram}). That $w_{\mathrm L}=w_{\mathrm H}-2$ is
part~(i) read for writhes and not a forward citation: switching $q$ makes its
sign $-1$ and lowers the writhe by two, and the oriented Reidemeister-II move
that follows removes $q$ and $y$, whose signs $-1$ and $+1$ cancel; what is
left is $D_{\mathrm L}$. At $\epsilon=0$ the crossing $y$ is a self-crossing
of the first component, so the mixed crossings of $D_{\mathrm A}$ are the old
mixed ones and their signed count is $2\ell$; then
$w_{\mathrm H}=w_1+w_2+2\ell+2$ and $w_{\mathrm L}=w_1+w_2+2\ell$. At
$\epsilon=1$ the crossing $y$ is itself mixed, so the old mixed crossings count
$2\ell-1$; then $w_{\mathrm H}=w_1+w_2+(2\ell-1)+2$ and
$w_{\mathrm L}=w_1+w_2+2\ell-1$. Both branches are
$w_{\mathrm L}=w_1+w_2+2\ell-\epsilon$. No crossing is \emph{deleted} in this
count, and an earlier revision said the mixed ones were; they are retained in
$D_{\mathrm A}$ and enter through $\ell$.


\emph{(vi).} First suppose $\epsilon=1$, or $\epsilon=0$ with $T$ same-piece.
By~(i), switching the positive $q$ lowers the writhe by two; the subsequent
oriented Reidemeister-II deletion removes crossings of signs $-1$ and $+1$,
so its net writhe change is zero. Thus $w_{\mathrm H}=w_{\mathrm L}+2$.
If $\epsilon=0$ and $T$ is different-piece, Lemma~\ref{lem:chambertransport}, clause~(B)(iii)
gives the same identity. Corollary~\ref{cor:groupedknot}(B), including its stated
empty-carrier case, identifies each grouped carrier writhe $w_{T,C_\nu}$ with
the writhe $w(D_\nu)$ of its grouped diagram, while
Definition~\ref{def:markeddata} binds $w_\nu=w(D_\nu)$. Thus
$w_{T,C_{\mathrm H}}=w_{\mathrm H}$ and $w_{T,C_{\mathrm L}}=w_{\mathrm L}$.
In every case Lemma~\ref{lem:chambercommon}(ii) gives $\rot_{\mathrm H}=\rot_{\mathrm L}$, hence $R_{\mathrm H}=R_{\mathrm L}$. Finally,
\begin{align*}
 d_{\mathrm H}&=1-w_{\mathrm H}-R_{\mathrm H}\\
 &=1-(w_{\mathrm L}+2)-R_{\mathrm L}\\
 &=(1-w_{\mathrm L}-R_{\mathrm L})-2\\
 &=d_{\mathrm L}-2.
\end{align*}

\emph{(vii).} $W_T\neq0$ and Lemma~\ref{lem:selectorid}, clause~(B) at $\epsilon=1$ give
$W_T=-s_{\mathrm{low}}W_1W_2$, with $s_{\mathrm{low}}$ as bound in
Lemma~\ref{lem:selectorid}, clause~(B), so both half selectors are nonzero and both half
carriers $L_1,L_2$ are uniform (Definition~\ref{def:wind}). That is the route
the hypothesis actually supports: it is a statement about the full selector, and
the halves' uniformity is deduced from it. The spectator product $U_T$ is not
used anywhere in this proof, which is why the statement asks only for
$W_T\neq0$. Each is a closed polygon whose principal turns are all
nonzero and, being turns of a generic polygon, of magnitude below $\pi$; by
Corollary~\ref{cor:groupedknot}(B) the grouped object on $L_i$ is a single knot
diagram with all crossings positive, writhe $w_i$, underlying plane curve
$L_i$ and HOMFLY polynomial $Q_i$. Theorem~\ref{thm:carrierfloor}(C) applies to it
in its uniform case --- or Theorem~\ref{thm:carrierfloor}(D), if the half bears no
residual piece at all --- and gives
\[
   \min\deg_a Q_i\;\geq\;1-w_i-R_i=d_i,
\]
hence $\min\deg_af_i\geq d_i$ as well, $f_i$ being a coefficient row of $Q_i$.

Therefore $f_1f_2$ has no support below $D=d_1+d_2$, so
$[a^{D-2}]f_1f_2=0$; and at degree $D$ the only way to write $D$ as a sum of an
exponent of $f_1$ at least $d_1$ and one of $f_2$ at least $d_2$ is
$(d_1,d_2)$, so $[a^{D}]f_1f_2=\omega_1\omega_2$. The bracket is
$0-\omega_1\omega_2+\omega_1\omega_2=0$. If one of the rows is zero then $f_1f_2=0$ and $\omega_1\omega_2=0$, so all
three terms of the bracket are zero.

The restriction to $W_T\neq0$ is not cosmetic. A summand with $W_T=0$ carries no
uniform halves for Theorem~\ref{thm:carrierfloor}(C) to be applied to, and nothing
is asserted about its bracket; it contributes zero to the weighted sum because
its weight is zero. An earlier revision wrote $W_TU_T\neq0$ here and in the
opening line while stating the proposition for $W_T\neq0$; the two stale
occurrences are corrected, and the hypothesis is the weaker one throughout.

\emph{(viii).} Put $L^*=C_{\mathrm H}$ and $q:=x$. Let $D_{\mathrm A}$ be the oriented
smoothing of $D_{\mathrm H}$ at $q$ and put $F_{\mathrm A}=P_{D_{\mathrm A}}$.
By clause~(ii), $D_{\mathrm A}$ has exactly two ordered
components with rows $f_1,f_2$; let $\ell$ be their linking number.

Clause~(i) makes $q$ positive and supplies (T1) and (T2);
clause (vi) gives $d_{\mathrm H}=d_{\mathrm L}-2$ for
Lemma~\ref{lem:fulltwist}, and clause (iv) supplies the auxiliary row.
Corollary~\ref{cor:groupedknot}(B), including its stated empty-carrier case,
identifies the grouped chamber diagrams' underlying curves as
$C_{\mathrm L},C_{\mathrm H}$, so their marked and abstract $d_\nu,F_\nu,\Omega_\nu$ bindings agree. Therefore
\begin{align}
 \Omega_{\mathrm H}-\Omega_{\mathrm L}
 &=\bigl[a^{d_{\mathrm L}-1}z^{-1}\bigr]F_{\mathrm A}\notag\\
 &=\bigl[a^{d_{\mathrm L}+2\ell-2}\bigr]f_1f_2
   -\bigl[a^{d_{\mathrm L}+2\ell}\bigr]f_1f_2. \tag{FT1}
\end{align}
At $\epsilon=1$, Lemmas~\ref{lem:rotshift} and~\ref{lem:chambercommon}(ii),
together with clause~(v), give
\begin{align}
 \rot_{\mathrm L}&=\rot(C_{\mathrm L})=\rot(C_{\mathrm H})
 =\rot(L^*)=\rot_1+\rot_2,\notag\\
 2\ell&=w_{\mathrm L}+1-w_1-w_2. \tag{FT2}
\end{align}
Put $\Delta_R=R_1+R_2-R_{\mathrm L}$. The first line of (FT2) and the triangle
inequality give $\Delta_R\geq0$, while the definitions and its second line give
\begin{align}
 D&=2-w_1-w_2-R_1-R_2\notag\\
  &=2-(w_{\mathrm L}+1-2\ell)-(R_{\mathrm L}+\Delta_R)\notag\\
  &=d_{\mathrm L}+2\ell-\Delta_R. \tag{FT3}
\end{align}
Thus (FT1) becomes
\begin{equation}
 \Omega_{\mathrm H}-\Omega_{\mathrm L}
 =\bigl[a^{D+\Delta_R-2}\bigr]f_1f_2
  -\bigl[a^{D+\Delta_R}\bigr]f_1f_2. \tag{FT4}
\end{equation}

Since $W_T\neq0$, Definition~\ref{def:wind} makes the contact carrier
$L^*=C_{\mathrm H}$ uniform. The uniformity clause of
Lemma~\ref{lem:selectorid}, clause~(B) at $\epsilon=1$ then makes both half carriers
uniform with the same turn sign. Lemma~\ref{lem:uniformrot}(i) gives
$\rot_1,\rot_2$ that same sign, so
$|\rot_1|+|\rot_2|=|\rot_1+\rot_2|$. By (FT2) this says
\begin{equation}
   \Delta_R=R_1+R_2-R_{\mathrm L}=0. \tag{FT5}
\end{equation}
With (FT5), the right side of (FT4) is exactly the first two terms of
clause~(vii)'s displayed bracket. That exported bracket gives
$(\Omega_{\mathrm H}-\Omega_{\mathrm L})+\omega_1\omega_2=0$, which proves
the claim.
\end{proof}

\begin{remark}[attribution]\label{rem:fulltwistsource}
The counterpart's full-twist calculation is retained in the proof of the live
jump: (FT1) is the abstract skein and auxiliary two-component row; (FT2) is the
signed-rotation and writhe ledger; (FT3)--(FT4) are the rotation-defect shift;
and (FT5) is the live-support collapse. The final sentence uses the exported
combined bracket of Lemma~\ref{lem:triplebridge}(vii). The writhe identity
$2\ell=w_{\mathrm L}+1-w_1-w_2$ is also the ledger that refutes the withdrawn
route of Remark~\ref{rem:withdrawnbracket}.
\end{remark}

\subsection{The noninterlacing join}\label{sec:nijoin}

The calculation of §\ref{sec:fulltwist} does transfer to $\epsilon=0$, but not
by copying: one step changes and it changes two of the bookkeeping identities
with it. Throughout, $F_{\mathrm H}$ is the grouped polynomial of the contact
carrier in the two-crossing chamber and $F_{\mathrm L}$ in the zero-crossing
chamber. The symbols $F_{\mathrm A}$ and $\ell$ are \emph{not} bound here:
"the link obtained by smoothing one newborn" does not choose the newborn, and
Lemma~\ref{lem:triplebridge}'s clause~(iii) is not equivariant under exchanging
the two, so the choice is a binding and it is made inside the lemma that reads
it. The chamber data are
$w_{\mathrm L}$, $R_{\mathrm L}$ and $d_{\mathrm L}=1-w_{\mathrm L}-R_{\mathrm
L}$ for the low chamber and likewise with $\mathrm H$; the half data are
$w_i,R_i,d_i=1-w_i-R_i$ and $D=d_1+d_2$, as in
Definition~\ref{def:markeddata}. An earlier revision wrote
$d_{\mathrm L}=1-w_{\mathrm F}-R_{\mathrm L}$, mixing a symbol from a naming
scheme this section no longer uses.

\begin{lemma}[the noninterlacing twist, its ledgers and live rotation defect]\label{lem:nitransport}
\begin{enumerate}
\item[(A)] Here eligibility is as in Definition~\ref{def:eligible}.
Let $\epsilon=0$ and let $T$ be eligible \emph{and same-piece}: after smoothing
$T$ the two newborns lie in one residual piece and on one carrier. Bind here, not by reference: let $q:=x$ be the newborn whose first visit opens
the contact word of Lemma~\ref{lem:contactword} --- the smoothed crossing is
this one and not the other newborn --- let $D_{\mathrm A}$ be the oriented
smoothing of $D_{\mathrm H}$ at $q$ and $F_{\mathrm A}=P_{D_{\mathrm A}}$; by
Lemma~\ref{lem:triplebridge}(ii) at $\epsilon=0$, $D_{\mathrm A}$ has exactly
two components, ordered by the half-label convention, and $\ell$ is their
linking number. The ledgers below are the ones clauses (ii)--(iv) of that
lemma supply. Let $s\in\{\pm1\}$ be the sign
\eqref{eq:S6} in the initial chamber; the
configuration with $s=-1$ is the reflection of the one with $s=+1$, and every
determinant below negates with it, so the statement is written for $s=+1$ and
read for $s=-1$ through that reflection. Then:
\[
   F_{\mathrm H}=a^{-2}F_{\mathrm L}+a^{-1}zF_{\mathrm A},
   \qquad
   \Omega_{\mathrm H}-\Omega_{\mathrm L}=\bigl[a^{d_{\mathrm L}-1}z^{-1}\bigr]F_{\mathrm A},
\]
and, with $f_i,d_i,D$ as in Definition~\ref{def:markeddata} and
$\Delta_R=R_1+R_2-R_{\mathrm L}$,
\[
   2\ell=w_{\mathrm L}-w_1-w_2,
   \qquad
   D=d_{\mathrm L}+2\ell+1-\Delta_R,
   \qquad
   \Omega_{\mathrm H}-\Omega_{\mathrm L}
   =\bigl[a^{D+\Delta_R-3}\bigr]f_1f_2
    -\bigl[a^{D+\Delta_R-1}\bigr]f_1f_2 .
\]
The display is the $\epsilon=0$ specialization, this lemma's own hypothesis ---
an earlier revision printed the unified two-branch exponents
$D+\Delta_R+\epsilon-3$, $D+\Delta_R+\epsilon-1$ over the ledger
$D=d_{\mathrm L}+2\ell+1-\epsilon-\Delta_R$ and called them "valid at both
branches" inside a statement that fixes $\epsilon=0$; the unified form is the
peer's blueprint's and lives in the withdrawn-obstruction remark at the end of
this section, where both branches are in scope. The value of
$\Delta_R$ on a live noninterlacing support is supplied separately, and this
statement carries the exponents in $\Delta_R$ alone rather than the numbers
they become.
\item[(B)] Here eligibility is as in Definition~\ref{def:eligible}.
Let $T$ be eligible at a simple bigon event with $\epsilon=0$, and suppose
$W_T\neq0$ --- the selector of Definition~\ref{def:markeddata}, one number for
both chambers; an earlier revision wrote "the contact carrier's selector",
which names no bound quantity, and left $T$ unbound. Write
\[
   \Delta_R=R_1+R_2-R_{\mathrm L}
\]
for the rotation defect, in the absolute rotations $R_1,R_2$ of the two half
contact carriers and $R_{\mathrm L}$ of the contact carrier in the low chamber, all bound in
Definition~\ref{def:markeddata}. Put $D=d_1+d_2$, with the half slots $d_i$
bound there. Then $\Delta_R=-1$; substituted into the two exponents of clause~(A),
that gives $D-4$ and $D-2$. The symbol $\Delta_R$ is rebound here: clause~(B)
has its own hypotheses and does not inherit clause~(A)'s local binding.
\end{enumerate}
\end{lemma}
\begin{proof}
\emph{Clause (A).} The first two displays are the two clauses of Lemma~\ref{lem:fulltwist}. Its
hypotheses (T1) and (T2) hold here by Lemma~\ref{lem:triplebridge}(i) --- that
lemma proves them at \emph{both} branches, so nothing is imported from
$\epsilon=1$ and the word ``verbatim'', which an earlier revision used here, is
gone. The \emph{second} clause carries a third antecedent,
$d_{\mathrm H}=d_{\mathrm L}-2$, which an earlier revision of this proof did not
name: it is Lemma~\ref{lem:triplebridge}(vi), stated for an eligible $T$ at a simple
bigon event and so available at this branch as well. Without it the coefficient
display is not licensed, the first display being unconditional and the second
not.

The step that differs between the branches is the identification of the two
components of the auxiliary link, and it too is
Lemma~\ref{lem:triplebridge}(ii)--(iii). At $\epsilon=1$ the chords interlace,
smoothing $x$ puts the two visits of $y$ on different components, and $y$ is a
mixed crossing --- \emph{retained}, in neither component's own diagram, and
counted in the linking exponent, which is the typing of
Lemma~\ref{lem:triplebridge}(ii); an earlier revision here said it was to be
deleted. At $\epsilon=0$ the chords are nested: smoothing
$x$ cuts the
traversal into $y_0Ay_1$ and $B$, so both visits of $y$ land on the first
component and $y$ survives there --- as a \emph{positive self Reidemeister-I
curl}. Deleting that curl changes no polynomial (Axiom~\ref{ax:homfly}) and
leaves exactly the two half groupings, so Lemma~\ref{lem:homflyrows}(ii) applies with rows
$f_1,f_2$ and gives $[z^{-1}]F_{\mathrm A}=a^{-2\ell}(a-a^{-1})f_1f_2$.

That deletion is also what changes the writhe ledger. Smoothing one positive
newborn removes one crossing from the two-crossing diagram and deleting the
curl removes one more, so $2\ell=w_{\mathrm L}-w_1-w_2$ --- without the $+1$
that Lemma~\ref{lem:triplebridge}(v) carries at $\epsilon=1$. Then
\[
   D=2-w_1-w_2-R_1-R_2
    =2-(w_{\mathrm L}-2\ell)-(R_{\mathrm L}+\Delta_R)
    =d_{\mathrm L}+2\ell+1-\Delta_R,
\]
and substituting into
$\bigl[a^{d_{\mathrm L}-1}\bigr]\bigl(a^{-2\ell}(a-a^{-1})f_1f_2\bigr)
=[a^{d_{\mathrm L}+2\ell-2}]f_1f_2-[a^{d_{\mathrm L}+2\ell}]f_1f_2$
with $d_{\mathrm L}+2\ell=D+\Delta_R-1$ gives the last display.

\emph{Clause (B).} The signed rotations $\rot_1=\rot(L_1)$ and $\rot_2=\rot(L_2)$ of the two half
contact carriers are bound in Definition~\ref{def:markeddata}, together with
$R_i=|\rot_i|$; write $\rot_{\mathrm L}$ for the signed rotation of the contact
carrier in the low chamber, so that $R_{\mathrm L}=|\rot_{\mathrm L}|$. An
earlier revision rebound these as $r_1,r_2$ here and as $\rho_i$ in the
definition, three names for one quantity. Let $s_{\mathrm{low}}$ be
the low-chamber side sign of Lemma~\ref{lem:rotshift} and $\tau$ the sign of the
turn at $M$ on the contact carrier. In the chart of Lemma~\ref{lem:contactturnboth}, normalised by
$s_{\mathrm{low}}=+1$, the two half contact turns are $\arg q$ and
$\pi-\arg p$, both in $(0,\pi)$; so the common sign of the two half contact
turns is $s_{\mathrm{low}}$, both sides negating together under the reflection
that sends $s_{\mathrm{low}}$ to $-1$. By Lemma~\ref{lem:contactturnboth} at
$\epsilon=0$ that sign and $\tau$ are opposite, $\tau=-s_{\mathrm{low}}$. All
noncontact corners partition between the halves, and
Lemma~\ref{lem:rotshift} at $\epsilon=0$ says the two halves carry one full
turn more than the carrier, of sign $s_{\mathrm{low}}$:
$\rot_1+\rot_2-\rot_{\mathrm L}=s_{\mathrm{low}}=-\tau$.

Suppose the selector is nonzero. Then the contact carrier is uniform
(Definition~\ref{def:wind}), so Lemma~\ref{lem:uniformrot}(i) gives
$\tau\rot_{\mathrm L}\geq1$. Each half inherits only noncontact turns of sign
$\tau$, together with exactly one contact turn of the opposite sign and of
magnitude strictly below $\pi$ (Lemma~\ref{lem:contactturnboth}(iv)). Apply
Lemma~\ref{lem:uniformrot}(ii) directly when $\tau=+1$ and after orientation
reversal when $\tau=-1$; it gives
$\tau\rot_1\geq1$ and $\tau\rot_2\geq1$. The absolute values may therefore all be
removed with the common sign $\tau$:
\[
   \Delta_R=|\rot_1|+|\rot_2|-|\rot_{\mathrm L}|
   =\tau\bigl(\rot_1+\rot_2-\rot_{\mathrm L}\bigr)=\tau(-\tau)=-1 .
\]
\end{proof}

\begin{theorem}[target factorization, dominated and newborn rows, both bigon branches and universal (S7)]\label{thm:s7universal}
If $T\in\Ind(G_{P_0})$ meets $\mathcal N$ then
$N(T)-L(T)+X(T)+Y(T)=0$.
\begin{enumerate}
\item[(A)] Here eligibility is as in Definition~\ref{def:eligible}.
Let $T$ be eligible with parts $T_A,T_B$, and let $W_1,W_2,U_T,\omega_1,\omega_2$
be as in Definition~\ref{def:markeddata}. Then
\[
   W_1W_2\,U_T\,\omega_1\omega_2
   \;=\;
   \Bigl(\wind_{\lambda_1}(T_A)\prod_{L}\Omega_1(T_A,L)\Bigr)
   \Bigl(\wind_{\lambda_2}(T_B)\prod_{L'}\Omega_1(T_B,L')\Bigr),
\]
the first product running over the carriers of $T_A$ in $\lambda_1$ and the
second over the carriers of $T_B$ in $\lambda_2$, as in
Definition~\ref{def:X1};
so that the two factors on the right are the terms of $T_A$ in
$X_1(\lambda_1)$ and of $T_B$ in $X_1(\lambda_2)$.
\item[(B)] Here eligibility is as in Definition~\ref{def:eligible}, and all
marked-data symbols are as in Definition~\ref{def:markeddata}.
Let the event be a simple vertex-on-edge event of bigon type, in the sense of
Convention~\ref{conv:events}, whose two newborn crossings interlace; hence
$\epsilon=1$. Then:
\begin{enumerate}
\item[(i)] For every eligible $T$, write
\[
   W_x=\wind_{P_2}\bigl(T\cup\{x\}\bigr),\qquad
   W_y=\wind_{P_2}\bigl(T\cup\{y\}\bigr)
\]
for the two one-newborn selectors. Then $W_x=W_y=0$, so
$X(T)=Y(T)=0$, whether or not the contact carrier of $T$ is uniform.
\item[(ii)] For every eligible $T$, define its returned contribution and
target share by
\[
   R_T=\delta\bigl(N(T)-L(T)+X(T)+Y(T)\bigr),\qquad
   J_T=s_{\mathrm{init}}\,W_1W_2\,U_T\,\omega_1\omega_2,
\]
where $\delta=\pm1$ is the wall orientation of Definition~\ref{def:branches},
$J$ is the sector target of Definition~\ref{def:sectors} --- an earlier revision
attributed it to Definition~\ref{def:branches} --- and
$s_{\mathrm{init}}=\delta\,s_{\mathrm{low}}$ is the \eqref{eq:S6} sign in the
initial chamber. Then $R_T=J_T$ when $W_T\neq0$; when
$W_T=0$, one has $J_T=0$ and $R_T=\delta(X(T)+Y(T))$; and
\[
   F=J\ \text{ at the interlacing bigon}
   \quad\Longleftrightarrow\quad
   \sum_{\substack{T\text{ eligible}\\ W_T=0}}
   \bigl(X(T)+Y(T)\bigr)=0 .
\]
\item[(iii)] If the two halves are generic polygons --- the hypothesis (S7)
carried in Definition~\ref{def:laws}, discharged at every such event on a
path in $\mathcal R_n$ by the certificate of the induction --- then $X_1$
satisfies (S7) there, with the sign \eqref{eq:S6} and the halves of
Definition~\ref{def:halves}.
\end{enumerate}
\item[(C)] Here eligibility is as in Definition~\ref{def:eligible} and the marked-data symbols are
as in Definition~\ref{def:markeddata}.
Let $\epsilon=0$. For every eligible $T$ with $W_T\neq0$ --- the selector alone,
the spectator product $U_T$ playing no part ---
\[
   \bigl[a^{D-4}\bigr]f_1f_2=\bigl[a^{D-2}\bigr]f_1f_2=0 .
\]
\item[(D)] The following local statements hold. In clauses~(ii)--(v),
eligibility is as in Definition~\ref{def:eligible}; let $\epsilon=0$ and let
$T$ be eligible.
\begin{enumerate}
\item[(i)] Let $S\in\Ind(G_P)$, let $A$ be a uniform carrier of $S$, and let $\{c\}$ be a
singleton residual piece of $S$ carried by $A$. Then
\[
   \min\deg_af_A\;\geq\;\bigl(1-w_{S,A}-R(A)\bigr)+2,
\]
so the factor $\Omega_1(S,A)$ of Definition~\ref{def:X1} is zero.
\item[(ii)] The set $T\cup\{x\}$ is independent and $\{y\}$ is a singleton
residual piece of it; symmetrically, $T\cup\{y\}$ is independent and $\{x\}$
is a singleton residual piece of it. No hypothesis on the pieces of $T$ is
used.
\item[(iii)] $X(T)=Y(T)=0$.
\item[(iv)] If $T$ is different-piece, put $S=T\cup\{x,y\}$. Then:
\begin{enumerate}
\item[(a)] $S$ is independent and $U(S)=U(T)\setminus\{x,y\}$;
\item[(b)] the residual pieces of $S$ are exactly the residual pieces of $T$
other than the two singletons $\{x\}$ and $\{y\}$ --- the same label sets, hence
the same piece diagrams, polynomials and writhes;
\item[(c)] with $L_1,L_2$ the two half contact carriers of
Lemma~\ref{lem:contactsplit}, rows $f_1,f_2$ and grouped writhes $w_1,w_2$, the
high-neither contact carrier satisfies $f_{C_{\mathrm H}}=f_1f_2$ and
$w_{T,C_{\mathrm H}}=w_1+w_2+2$, where $C_{\mathrm H}$ denotes that carrier;
\item[(d)] consequently $f_{C_{\mathrm L}}=f_1f_2$ and $w_{T,C_{\mathrm L}}=w_1+w_2$ in the low
chamber.
\end{enumerate}
\item[(v)] If $T$ is different-piece, the entire row vanishes:
\[
   N(T)=L(T)=X(T)=Y(T)=0 .
\]
\end{enumerate}
\item[(E)] Let the event be a simple vertex-on-edge event of bigon type, in the sense of
Convention~\ref{conv:events}, whose two newborn crossings do not interlace, and
whose two halves are generic polygons. Then $X_1$ satisfies (S7) there, with the
sign \eqref{eq:S6} and the halves of Definition~\ref{def:halves}.
\item[(F)] Let $W$ be a simple vertex-on-edge event in the sense of
Convention~\ref{conv:events}, of either the bigon or the sliding type, with any
number of exterior strands threading the region of the event, and suppose the
two halves of Definition~\ref{def:halves} are generic polygons with fewer
vertices than the parent. Then $X_1$ satisfies (S7) at $W$ with the sign (S6)
and the halves (S1) of Definition~\ref{def:halves}.

The hypothesis on the halves is not a restriction: it holds at every simple
vertex-on-edge event on a path inside $\mathcal R_n$, which is the certificate
proved in the induction's section. The genericity and vertex-count hypotheses are stated here for different reasons, and an earlier
revision gave one reason for both. The genericity is \emph{consumed} by this
proof, which evaluates $X_1$ on each half and cannot do so otherwise. The
vertex count is \emph{exported}: no step below reads it --- this theorem is not
an induction and applies no inductive hypothesis --- and it is carried because
the consumer is, the induction of the last section applying its hypothesis to
the halves and needing them smaller than the parent.
\end{enumerate}
\end{theorem}
\begin{proof}
\emph{The opening assertion.} Let $q\in T\cap\mathcal N$. Then $q$ is adjacent to both newborns
(Lemma~\ref{lem:contactword}(ii)), so neither $T\cup\{x\}$ nor $T\cup\{y\}$ is
independent and $X(T)=Y(T)=0$. Also $x,y\in N_{G_{P_2}}(T)$, so neither is in
$U(T)$: the newborns are dominated and are not residual crossings. Hence the
residual sets, the pieces, and their diagrams are computed from the old labels
alone and coincide in the two chambers, the crossing data being constant across
the event (Lemma~\ref{lem:guardconst}); the carriers are the cycles of the same
smoothing of the same word, with the same corner signs; and the rotations agree
because the two chambers are joined by a path along which no carrier turn
vanishes. So every factor of the term agrees and $N(T)=L(T)$.

\emph{Clause (A).} By Definition~\ref{def:X1} the term of $T_A$ is its selector times the product
of the reads over all of its carriers, and $\wind_{\lambda_1}(T_A)=W_1$ by
Definition~\ref{def:markeddata}; likewise for $T_B$. So the right-hand side is
$W_1W_2$ times the product of the reads over all carriers of $T_A$ in
$\lambda_1$ and of $T_B$ in $\lambda_2$. It remains to identify that product
with $U_T\omega_1\omega_2$, carrier by carrier.

\emph{The contact carriers.} $L_1$ and $L_2$ are the two halves' contact
carriers, and $\Omega_1(T_A,L_1)=\omega_1$, $\Omega_1(T_B,L_2)=\omega_2$ by
Definition~\ref{def:markeddata}, which binds $\omega_i$ as exactly that read.

\emph{The others.} Lemma~\ref{lem:contactsplit} carries the carriers of $T$ in
the two-crossing chamber other than $L^{*}$ bijectively to the carriers of
$T_A$ and $T_B$ other than $L_1$ and $L_2$. A carrier in that correspondence
bears the same residual pieces on both sides: no retained \emph{mixed} label
lies on it, since a mixed label has one visit in each interval and so its
carrier meets both intervals, which by Lemma~\ref{lem:contactsplit} is $L^{*}$
alone; and for the unmixed labels Lemma~\ref{lem:ownermap} identifies the
residual sets, $U(T)\cap V_A=U_{\lambda_1}(T_A)$ and
$U(T)\cap V_B=U_{\lambda_2}(T_B)$, while the induced interlacement graph on the
labels of one interval is the half's own. Equal label sets give equal piece \emph{records}: the map is the identity on the
retained visits, it preserves their cyclic order because the two carriers are
the same closed curve traversed the same way, and it preserves the over/under
designation and the sign at each visit, those being read from the two branch
directions (Definition~\ref{def:piecediagram}). That is a record isomorphism in
the sense of Definition~\ref{def:record}, so Axiom~\ref{ax:gausscode} gives
equal links and hence equal polynomials, and the writhes agree as crossing
counts. An earlier revision passed from equal label sets to equal polynomials
with no isomorphism and no axiom. The carrier being the same closed curve with
the same corners, the rotation agrees too; so the two reads agree, and
their product over the matched carriers is $U_T$, again by
Definition~\ref{def:markeddata}.

Multiplying the three displays gives the identity. An earlier revision asserted
it as part of the ordered-product bijection, which matches \emph{supports} and
says nothing about owners, rows, writhes or rotations.

\emph{Clause (B).} \emph{Clause (i).}
\emph{The collar invariant.} By eligibility and Lemma~\ref{lem:contactword},
both visits of every selected old label lie wholly in $A$ or wholly in $B$. So
no smoothing of $T$ cuts the boundary collar between a newborn smoothing corner
and $M$: in the $x$-branch the smoothing corner at the remote visit $x_1$
reconnects to the position immediately after the incident visit $x_0$, and that
uninterrupted arc reaches $M$ before any selected visit; in the $y$-branch the
uninterrupted incident arc from $M$ reaches $y_0$ before any selected visit and
then closes at the remote visit $y_1$. Hence, for every eligible $T$ including
nonempty ones, the relevant newborn smoothing corner and the turn at $M$ lie on
one and the same daughter carrier.

\emph{The signs, computed in the high chamber.} Work where the two newborns
actually exist. Put $M=(0,0)$, let the remote direction be $r=(1,0)$ on the line
$y=c$ with $c>0$, and write $p-M=(u,h)$ and $q-M=(v,k)$ with $h,k>c$, so that
both neighbours lie beyond the remote line. Interlacing is $u/h<v/k$, i.e.\
$hv-uk>0$. The three determinants that matter are then
\[
   \det\bigl((-u,-h),(v,k)\bigr)=hv-uk>0,
   \qquad
   \det\bigl(r,(-u,-h)\bigr)=-h<0,
   \qquad
   \det\bigl((v,k),r\bigr)=-k<0 .
\]
The first is the turn at $M$, from the incoming direction $-(u,h)$ to the
outgoing direction $(v,k)$: a \emph{left} turn. The second is the corner of the
daughter carrier produced by smoothing $x$, which arrives along $r$ and leaves
along $-(u,h)$: a \emph{right} turn. The third is the corner of the daughter
produced by smoothing $y$, which arrives along $(v,k)$ and leaves along $r$:
again a right turn. By the collar invariant each of those corners shares its
carrier with the turn at $M$, so each of those carriers is mixed, its weight is
$0$ by Definition~\ref{def:wind}, and $W_x=W_y=0$. Reflecting the configuration
negates all three determinants together and leaves every carrier mixed.


\emph{Clause (ii).}
$J_T$ has a second form, used throughout below:
$W_T=-s_{\mathrm{low}}W_1W_2$ by Lemma~\ref{lem:selectorid}, clause~(B), so
$J_T=-\delta\,W_TU_T\,\omega_1\omega_2$. The statement carries the first form
only, that being the one built from the two half terms; an earlier revision put
the derivation inside the statement.

Both quantities are directed: each carries the wall orientation $\delta$, so the
equivalence compares like with like. An earlier revision wrote an undirected row
$N-L-J(T)$ with $J(T)$ never defined, against Definition~\ref{def:branches}'s
directed $R=\delta\sum(N-L+X+Y)$; the display above is the typed replacement.

\emph{If $W_T\neq0$.} By clause~(B)(i) both one-newborn
branches are dead, so $X(T)=Y(T)=0$ and $R_T=\delta(N(T)-L(T))$. Every carrier
other than the contact carrier is common to the two chambers with the same data,
and the selector is the same in both, by
Lemma~\ref{lem:chambercommon}, clauses (i) and (ii); with $W_T$ and $U_T$ bound separately as in
Definition~\ref{def:markeddata},
\[
   N(T)-L(T)=W_TU_T\bigl(\Omega_{\mathrm H}-\Omega_{\mathrm L}\bigr)
   =-W_TU_T\,\omega_1\omega_2
\]
by Lemma~\ref{lem:triplebridge}(viii). So
$R_T=-\delta W_TU_T\omega_1\omega_2=J_T$.

\emph{If $W_T=0$.} Then $J_T=0$ by its own display, and the high-neither and low
terms vanish with their common prefactor
(Lemma~\ref{lem:chambercommon}, clauses (i) and (ii), which are what the
prefactor is made of), leaving $R_T=\delta(X(T)+Y(T))$.

\emph{Summation.} By the opening assertion, supports meeting $\mathcal N$
contribute nothing. That $\sum_TJ_T$ is the global target $J$ needs two
statements and an earlier revision cited only the first.
Lemma~\ref{lem:orderedproduct} matches the eligible supports bijectively with
the pairs $(T_A,T_B)$; clause~(A) matches each product
$W_1W_2U_T\omega_1\omega_2$ with the product of the two half terms. Summing the
second over the bijection of the first gives
$\sum_TW_1W_2U_T\omega_1\omega_2=X_1(\lambda_1)X_1(\lambda_2)$, and
$J=s_{\mathrm{init}}X_1(\lambda_1)X_1(\lambda_2)$ by
Definition~\ref{def:sectors}, which is where the sector target is defined --- the
statement above says so and this proof said otherwise, the statement having been
repaired in an earlier revision and the proof left alone. Only after that identification is the common
nonzero $\delta$ cancelled from $\sum R_T-\sum J_T$, which gives the
equivalence.


\emph{Clause (iii).}
By Proposition~\ref{prop:sectortarget}(F) it suffices that $R=\epsilon J=J$, and by clause~(B)(ii) that is
$\sum_{T\ \mathrm{eligible},\,W_T=0}\bigl(X(T)+Y(T)\bigr)=0$. By clause~(B)(i) every term of that sum is zero.

\emph{Clause (C).} \emph{The two forbidden degrees.} $W_T\neq0$ means the full carrier through
the contact is live, so by Definition~\ref{def:wind} it turns uniformly. By
Lemma~\ref{lem:contactturnboth} in the branch $\epsilon=0$ its contact turn has
the sign of $-s_{\mathrm{low}}$ while each half's contact turn has the sign of
$s_{\mathrm{low}}$ --- the low-chamber sign of Lemma~\ref{lem:selectorid}, clause~(B), which
is what Lemma~\ref{lem:contactturnboth} states and what an earlier revision here
wrote as a bare $s$; and by
Lemma~\ref{lem:contactsplit} each half carrier carries the full carrier's
non-contact corners together with its own contact turn. So each half carrier
has all its turns of one sign except exactly one, of the opposite sign and of
magnitude strictly below $\pi$ (Lemma~\ref{lem:contactturnboth}(iv)). That is the
one-dissent hypothesis of Theorem~\ref{thm:carrierfloor}(C), which applies by
Corollary~\ref{cor:groupedknot}(B) to the grouped knot on each half, or
Theorem~\ref{thm:carrierfloor}(D) if that half bears no piece, and gives
$\min\deg_af_i\geq d_i$. Hence $f_1f_2$ has no support below $D$, and in
particular none at $D-4$ or $D-2$.


\emph{Clause (D).} \emph{Clause (i).} \emph{Pass to the larger support.} Put $S'=S\cup\{c\}$. It is independent: $c$
is residual, so it is adjacent to no member of $S$. Its undominated set is
\[
   U(S')=U(S)\setminus\bigl(\{c\}\cup N_{G_P}(c)\bigr)=U(S)\setminus\{c\},
\]
because $\{c\}$ is a component of $G_P[U(S)]$, so $c$ has no neighbour inside
$U(S)$. Deleting an isolated vertex leaves every other component of a graph
unchanged, so the residual pieces of $S'$ are exactly the residual pieces of $S$
other than $\{c\}$ --- the same label sets, hence by
Definition~\ref{def:piecediagram} the same piece diagrams, the same $P_H$ and
the same $|H|$. Nothing is being compared across a change of diagram here: the
pieces are literally the same objects.

\emph{The two loops are carriers.} Smoothing $c$ cuts the traversal of $A$ at
the two visits of $c$ and reconnects, so the carriers of $S'$ are those of $S$
with $A$ replaced by the two closed curves $\Lambda_1,\Lambda_2$ carrying the
two arcs of $A$ between the visits of $c$; every other carrier of $S$ is
untouched, $c$ lying on $A$. Each residual piece of $S'$ formerly on $A$
therefore lies on $\Lambda_1$ or on $\Lambda_2$, by
Lemma~\ref{lem:carriers}(iv) applied to $S'$; no separate no-straddle
argument is needed and no sub-loop is compared with its parent. Hence
\[
   P_{S,A}=P_{S',\Lambda_1}\,P_{S',\Lambda_2}\cdot P_{\{c\}}
   =P_{S',\Lambda_1}\,P_{S',\Lambda_2},
   \qquad
   w_{S,A}=w_{S',\Lambda_1}+w_{S',\Lambda_2}+1,
\]
the singleton contributing $P_{\{c\}}=1$ and one unit of writhe
(Definition~\ref{def:piecediagram}: a closed curve with one crossing is an
unknot diagram). Each factor is a knot polynomial, hence lies in
$\mathbb Z[a^{\pm1},z^2]$ and has no negative power of $z$ ---
Axiom~\ref{ax:homfly}, whose knot-parity consequence this step reads and which
an earlier revision used without naming --- so the zeroth rows multiply:
$f_A=f_{\Lambda_1}f_{\Lambda_2}$.

\emph{Turns and rotations.} $A$ has no corner at $c$, which is neither a
polygon vertex nor a member of $S$; the two smoothing corners of $\Lambda_1$ and
$\Lambda_2$ at the site of $c$ receive turning angles that are negatives of
each other, so by Lemma~\ref{lem:turnlift}(ii)
$\rot(A)=\rot(\Lambda_1)+\rot(\Lambda_2)$. Their determinants are $\det(u,v)$
and $\det(v,u)$ for the two branch directions at $c$, hence opposite and
nonzero by (G5). Let $\sigma$ be the common sign of $A$'s turns; every corner of
either loop other than the two new ones is a corner of $A$ and has sign
$\sigma$. So one loop is uniform and the other has exactly one dissenting
corner, of magnitude below $\pi$. Lemma~\ref{lem:uniformrot}(i) applies to the
uniform loop. Apply Lemma~\ref{lem:uniformrot}(ii) to the dissenting loop
directly when $\sigma=+1$ and after orientation reversal when $\sigma=-1$.
Thus $\sigma\rot(\Lambda_i)\geq1$ for both $i$, so the two rotations carry the
sign $\sigma$ and $R(A)=R(\Lambda_1)+R(\Lambda_2)$.

\emph{The bound.} Theorem~\ref{thm:carrierfloor}(C) applies to each loop as a
carrier of $S'$ --- through Corollary~\ref{cor:groupedknot}(B) if it bears a piece
and Theorem~\ref{thm:carrierfloor}(D) if it does not, one loop uniform and the other
one-dissent --- so
$\min\deg_af_{\Lambda_i}\geq1-w_{S',\Lambda_i}-R(\Lambda_i)$ when the row is
nonzero, and if either row vanishes then $f_A=0$ and the claim is trivial.
Adding the two bounds and substituting the two displays above,
\[
   \min\deg_af_A\;\geq\;2-\bigl(w_{S,A}-1\bigr)-R(A)
   =\bigl(1-w_{S,A}-R(A)\bigr)+2 .
\]
The factor $\Omega_1(S,A)$ reads the coefficient at $1-w_{S,A}-R(A)$, two
degrees below the row's lowest nonvanishing degree, so it is zero.

\emph{Clause (ii).} $x$ is nonadjacent to $y$ because $\epsilon=0$, and nonadjacent to every element
of $T$ because $T$ is eligible; so $T\cup\{x\}$ is independent. For the same two
reasons $y$ is dominated by no element of $T\cup\{x\}$, so it is undominated and
lies in some residual piece. Let $c$ be an old crossing adjacent to $y$. By
Lemma~\ref{lem:contactword}(ii) $c$ is adjacent to $x$ as well, and $x$ now lies in
the support, so $c$ is dominated and is not residual. Hence $y$ has no residual
neighbour and $\{y\}$ is a component of the residual graph. The statement and
every step of this proof are invariant under interchanging the names $x$ and
$y$, which gives the symmetric assertion.

\emph{Clause (iii).} Take the branch that smooths $x$; the statement and the argument are invariant
under interchanging the two names, which gives the other branch. If
$\wind(T\cup\{x\})=0$ the term is zero by Definition~\ref{def:X1}. Otherwise
every carrier of $T\cup\{x\}$ is uniform. By clause~(D)(ii) $\{y\}$ is a singleton residual piece of
$T\cup\{x\}$; let $A$ be the carrier bearing it
(Lemma~\ref{lem:carriers}(iv)). Then $A$ is uniform and carries a singleton,
so $\Omega_1(T\cup\{x\},A)=0$ by clause~(D)(i), and the term, a
product with that factor, is zero.

\emph{Clause (iv).}
For clause~(D)(iv)(a), at $\epsilon=0$ the newborns do not interlace, so they are nonadjacent, and
each is nonadjacent to every member of $T$ because $T$ is eligible; hence $S$ is
independent. For the undominated set, let $c$ be an old crossing lying in
$U(T)$, and suppose $c\in N_{G_{P_2}}(x)$. By Lemma~\ref{lem:contactword}(ii) $c$ is
then a neighbour of $y$ as well, so $x,c,y$ is a path inside $G_{P_2}[U(T)]$
and the two newborns lie in one residual component of $T$, contradicting the
different-piece hypothesis. So $U(T)$ meets neither $N_{G_{P_2}}(x)$ nor $N_{G_{P_2}}(y)$, and
removing $x,y$ from $U(T)$ removes nothing else:
$U(S)=U(T)\setminus\{x,y\}$.

For clause~(D)(iv)(b), by Lemma~\ref{lem:chambertransport}, clause~(A) $\{x\}$ and $\{y\}$ are components of
$G_{P_2}[U(T)]$, that is, isolated vertices. Deleting isolated vertices leaves
every other component of a graph unchanged, so by clause~(D)(iv)(a) the components of
$G_{P_2}[U(S)]$ are exactly the components of $G_{P_2}[U(T)]$ other than those
two. A residual piece is a label set, and
Definition~\ref{def:piecediagram} builds its diagram from that label set and the
parent polygon alone, so the diagrams, the $P_H$ and the $|H|$ are the same
objects, not corresponding ones.

For clause~(D)(iv)(c), let $(T_A,T_B)$ be the ordered parts of $T$ from
Lemma~\ref{lem:orderedproduct}. Smoothing $x$ and $y$ affects only the carrier
of $T$ that contains them,
which is $C_{\mathrm H}$; every other carrier of $T$ is a carrier of $S$
unchanged, with unchanged pieces. By Proposition~\ref{prop:sectortarget}(A) the carriers of
$S$ in $P_2$ are the carriers of $T_A$ in $\lambda_1$, those of $T_B$ in
$\lambda_2$, and one local carrier $U_{\mathrm{loc}}$ in the contact disc; combining with
Lemma~\ref{lem:contactsplit}, the three of these into which $C_{\mathrm H}$
splits are the two half contact carriers $L_1,L_2$ and $U_{\mathrm{loc}}$. By
Proposition~\ref{prop:sectortarget}(B) $U_{\mathrm{loc}}$ bears no piece, and each of $L_1,L_2$ has
the same pieces, polynomials and grouped writhe as its mate in the half.

Now every residual piece of $T$ carried by $C_{\mathrm H}$ other than $\{x\}$
and $\{y\}$ is, by clause~(D)(iv)(b), a residual piece of $S$, and by
Lemma~\ref{lem:carriers}(iv) applied to $S$ it lies on exactly one carrier of
$S$; that carrier is one of the three just named, and not $U_{\mathrm{loc}}$. So those pieces
distribute between $L_1$ and $L_2$, and
\[
   P_{T,C_{\mathrm H}}
   =P_{\{x\}}\,P_{\{y\}}\prod_{H\text{ on }L_1}P_H\prod_{H\text{ on }L_2}P_H
   =P_{S,L_1}P_{S,L_2},
   \qquad
   w_{T,C_{\mathrm H}}=1+1+w_1+w_2,
\]
using $P_{\{x\}}=P_{\{y\}}=1$ and $|\{x\}|=|\{y\}|=1$ from
Lemma~\ref{lem:chambertransport}, clause~(A). Each factor is a knot polynomial, hence in
$\mathbb Z[a^{\pm1},z^2]$ with no negative power of $z$ by
Axiom~\ref{ax:homfly}, so the zeroth rows multiply.

Clause~(D)(iv)(d) follows from clause~(D)(iv)(c) together with Lemma~\ref{lem:chambertransport}, clause~(B)(iii), which gives
$f_{C_{\mathrm L}}=f_{C_{\mathrm H}}$ and $w_{T,C_{\mathrm L}}=w_{T,C_{\mathrm H}}-2$.


\emph{Clause (v).}
The one-newborn terms are clause~(D)(iii). For the other two,
Lemma~\ref{lem:chambertransport}, clause~(B) gives
$L(T)=W_TU_T[a^{d_0}]f_{C_{\mathrm L}}$ and $N(T)=W_TU_T[a^{d_0-2}]f_{C_{\mathrm L}}$ with no
hypothesis attached, so it suffices to treat the two cases of the common
prefactor.

If $W_TU_T=0$ both terms are zero as products with a zero factor.

So suppose $W_TU_T\neq0$. Then $W_T\neq0$, so every carrier of $T$ is uniform,
in particular the contact carrier. Let $L_1,L_2$ be the two half contact carriers of
Lemma~\ref{lem:contactsplit} --- the names that lemma and
Definition~\ref{def:markeddata} use; an earlier revision wrote
$L_2^{\mathrm h}$ here and then never used the symbol --- with rows $f_1,f_2$, writhes
$w_1,w_2$, rotations $R_1,R_2$, slots $d_i=1-w_i-R_i$ and $D=d_1+d_2$ as in
Definition~\ref{def:markeddata}. Three facts about them are already proved and
are used unchanged:
\begin{itemize}
\item $f_{C_{\mathrm L}}=f_1f_2$ and $w_{T,C_{\mathrm L}}=w_1+w_2$, by
clause~(D)(iv)(d), which is where the passage from the support
$T$ to the support $T\cup\{x,y\}$ --- the support that
clauses~(A) and~(B) of Proposition~\ref{prop:sectortarget} are about --- is
carried out and the old residual set is shown to survive it intact;
\item $R_1+R_2-R(C_{\mathrm L})=-1$, by Proposition~\ref{prop:sectortarget}(C) and
Lemma~\ref{lem:rotshift} at $\epsilon=0$, whose shift is the low-chamber side
sign $s_{\mathrm{low}}$ and whose absolute values may be taken with the common
sign because the contact carrier is live (the computation is that of
Lemma~\ref{lem:nitransport}, clause~(B));
\item $\min\deg_af_i\geq d_i$ for $i=1,2$: each half carrier keeps the uniform
non-contact turns of the live full carrier and acquires exactly one contact
turn of the opposite sign and of magnitude below $\pi$
(Lemma~\ref{lem:contactturnboth}(iv)), which is
the one-dissent hypothesis of Theorem~\ref{thm:carrierfloor}(C), applied through
Corollary~\ref{cor:groupedknot}(B) to the grouped knot on each half or through
Theorem~\ref{thm:carrierfloor}(D) if a half bears no piece.
\end{itemize}
Hence $f_{C_{\mathrm L}}=f_1f_2$ has no support below $D$, while
\[
   d_0=1-w_{T,C_{\mathrm L}}-R(C_{\mathrm L})=1-(w_1+w_2)-(R_1+R_2+1)=d_1+d_2-2=D-2 ,
\]
so $[a^{d_0}]f_{C_{\mathrm L}}=0$ and, two degrees lower still,
$[a^{d_0-2}]f_{C_{\mathrm L}}=[a^{D-4}]f_{C_{\mathrm L}}=0$. Both terms vanish.

\emph{Clause (E).} By Proposition~\ref{prop:sectortarget}(F) it suffices that $R=\epsilon J=0$, i.e.\
that $\sum_T\bigl(N(T)-L(T)+X(T)+Y(T)\bigr)=0$ in the notation of
Definition~\ref{def:branches}, the sum running over eligible $T$ by the
opening assertion. Fix such a $T$ and let $W$ be the
common selector of its high-neither and low supports, common by
Lemma~\ref{lem:chambercommon}(ii).

Split the eligible supports into the two sectors of
Lemma~\ref{lem:chambertransport}, clause~(A).

\emph{Different-piece supports.} Every one of the four terms of such a row is
zero by clause~(D)(v), with no hypothesis on the selector.

\emph{Same-piece supports.} These are the remaining ones, and the rest of this
proof treats them.

The one-newborn terms vanish in either case, by
clause~(D)(iii), which needs no hypothesis on any selector; so
only $N(T)-L(T)$ is at issue, and the split is on the prefactor.

\emph{If $U_T=0$.} Both $N(T)$ and $L(T)$ carry $W_TU_T$ as a factor, so both
vanish and the row is zero. The factorization is read off
Definition~\ref{def:X1} and clauses (i) and (ii) of
Lemma~\ref{lem:chambercommon} --- an earlier revision attributed those two
clauses to Lemma~\ref{lem:chambertransport}, clause~(B), which kept only (iii) when the two
were split --- and not off the latter's closing display: by
Definition~\ref{def:X1} the term of $T$ in either chamber is its selector times
the product of the reads over all of its carriers; by (i) the carriers other
than the contact carrier are the same closed curves with the same pieces in both
chambers, so their reads multiply to the same $U_T$ in both; and by (ii) the
selector is the same $W_T$ in both. What is left in each chamber is the read on
its own contact carrier. The closing display of
Lemma~\ref{lem:chambertransport}, clause~(B) is \emph{not} available here: it is stated
under the different-piece hypothesis of clause (iii), and this sector is the
same-piece one. An earlier revision of this proof cited it anyway, in both
same-piece branches. The case
is written out because the result invoked below is a statement about
coefficients, and multiplying a coefficient conclusion by a vanishing spectator
product is not the same as deducing one.

\emph{If $U_T\neq0$ and $W\neq0$}, by
Lemma~\ref{lem:nitransport}, clauses~(A) and~(B)
\[
   N(T)-L(T)=WU_T\Bigl([a^{D-4}]f_1f_2-[a^{D-2}]f_1f_2\Bigr),
\]
both of whose coefficients are zero by clause~(C), whose
hypothesis is $W\neq0$ and which is therefore invoked inside its own quantifier.

\emph{If $W=0$}, the high-neither and low terms both vanish with their common
prefactor, by the same reading of Definition~\ref{def:X1} with
Lemma~\ref{lem:chambercommon}(ii): the selector is one number $W=W_T$ in both
chambers and multiplies the whole term in each.

Every term is zero, so $R=0$ and $F=J$.

\emph{Clause (F).} The sliding type is Lemma~\ref{lem:slidingcoeff}, clause~(D). The bigon type splits by the
interlacement bit of Definition~\ref{def:sectors}: the interlacing case is
clause~(B)(iii), and the noninterlacing case is
clause~(E). Those three cases are exhaustive and mutually
exclusive: a simple vertex-on-edge event is of exactly one of the two types by
Lemma~\ref{lem:dictionary}, and a bigon event has $\epsilon\in\{0,1\}$.

No bound on the number of crossings, supports or residual pieces occurs in any
of the three arguments, and none of them inspects the region away from the
isolated contact, so arbitrary exterior crossings and arbitrary threading are
included.
\end{proof}

\begin{remark}[the superseded deletion argument]\label{rem:kinkdeletion}
Two earlier revisions --- the one-newborn half of
Theorem~\ref{thm:s7universal}, clause~(C) as first printed, and the different-piece argument
now in Theorem~\ref{thm:s7universal}(D)(v) --- argued by deleting the
surviving positive Reidemeister-I singleton from the affected carrier $A$ and
asserting that the cleaned carrier $C$ satisfies $R_A=R_C+1$, one unit less
absolute rotation; the second said explicitly that the argument did not use the
selector. Both halves of that are wrong. A crossing's sign does not determine
the sign of the change in Whitney index.

Theorem~\ref{thm:s7universal}(D)(i) replaces the deletion. It does not delete
anything: it cuts the carrier at the singleton and applies the floor to each of
the two loops, and it uses the selector, which is what supplies the uniform
turning that makes the two rotations add with one sign. Because
Theorem~\ref{thm:s7universal}(D)(ii) holds in both sectors --- the surviving
newborn is always isolated once the other is in the support, its old
neighbourhood being the other's --- the repair covers the same-piece sector as well, and
Theorem~\ref{thm:s7universal}, clause~(C) no longer carries a one-newborn half at all: it
states the two vanishing degrees, and the one-newborn terms are
Theorem~\ref{thm:s7universal}(D)(iii) wherever they are needed.
\end{remark}

\begin{remark}[what the different-piece sector cost, and what it did not]
\label{rem:diffcensus}
The proof of clause~(D)(v) above establishes vanishing in exactly the degrees $D-2$ and $D-4$
that Theorem~\ref{thm:s7universal}, clause~(C) kills in the same-piece sector, and by
the same floor on the same two halves. The different-piece sector is therefore
not a separate mechanism: what it needed was the typing --- that the two
newborns are singleton kinks (Lemma~\ref{lem:chambertransport}, clause~(A)), that the
chamber transport of the row and the prefactor holds with no selector
hypothesis (Lemma~\ref{lem:chambertransport}, clause~(B)), and that the one-newborn branch
dies by cutting at its surviving singleton rather than by deleting it
(Theorem~\ref{thm:s7universal}(D)(i)).

Two withdrawn attempts are recorded, each failing on an explicit
configuration printed here. The first asserted, as an axiom, that a
different-piece eligible support is selector-dead;
the peer refuted it with a guarded $24$-gon whose contact carrier is live,
with its data --- low word $[0,0,1,1]$, high word
$[0,2,1,1,3,3,2,0]$, edgeless interlacement, $\rot=3$ in both chambers, low slot
$-4$, high slot $-6$. The second proved the low read to vanish by
a three-loop decomposition of the contact carrier; that route is correct but it
needed a diagram identification between a proper sub-loop and its parent that
was not printed formally, and the half split above needs no such bridge, so the
route was withdrawn in favour of machinery already in the document.
\end{remark}

\begin{remark}[what this proof replaces, and whose the repair is]\label{rem:d61bridge}
An earlier revision argued these signs in the chart \emph{at the wall}, where
the two newborns coincide at $M$, and then applied the conclusion to the high
chamber where they do not. The limiting signs were right and the argument was
in the wrong chamber; and it assumed rather than printed the collar invariant
above. Both repairs --- the eligible-$T$ collar invariant and the exact
high-chamber determinants --- are the counterpart's.
\end{remark}

\begin{remark}[why the aggregate turned out to be termwise]\label{rem:aggregateclosed}
An earlier revision recorded this sum as an \emph{aggregate} identity over
selector-dead supports. It is not an aggregate identity: every term vanishes separately, because at an
interlacing bigon the turn at the moving vertex and the smoothing corner of
either newborn have \emph{opposite} signs, so whichever newborn is smoothed,
the carrier that inherits the moving vertex is mixed. The computation is three
lines in the chart of Lemma~\ref{lem:contactturnboth} and it needs neither a
pairing on supports nor a coefficient argument.

The same computation explains why the noninterlacing branch is different and
harder. There $\tau_M=\arg q-\arg p-\pi$ is a \emph{right} turn, of the same
sign as both smoothing corners, so the carriers can be uniform and the
one-newborn branches can be live --- which is exactly why they must be killed by
Theorem~\ref{thm:s7universal}(D)(iii), cutting the affected carrier at the newborn that
survives in the support, and not by any degree argument.
\end{remark}

\begin{remark}[a withdrawn obstruction]\label{rem:whynotni}
An earlier revision of this document printed a remark asserting that the
full-twist calculation \emph{cannot} transfer to $\epsilon=0$, on the ground
that smoothing one newborn leaves the other as a self-crossing so that the two
components are not the half groupings. The observation is right and the
conclusion drawn from it was wrong: the surviving crossing is a positive
Reidemeister-I curl, deleting it changes no polynomial, and what is left is
exactly the two half groupings. What the difference actually costs is one unit
in the writhe ledger --- $2\ell=w_{\mathrm L}-w_1-w_2$ here against
$2\ell=w_0+1-w_1-w_2$ there --- and, with the rotation defect, that is the
whole of why the two branches read the diagonals $D-4,D-2$ and $D-2,D$. In one
formula, over both branches: the ledger is
$D=d_{\mathrm L}+2\ell+1-\epsilon-\Delta_R$ and the jump reads
\[
   \Omega_{\mathrm H}-\Omega_{\mathrm L}
   =\bigl[a^{D+\Delta_R+\epsilon-3}\bigr]f_1f_2
    -\bigl[a^{D+\Delta_R+\epsilon-1}\bigr]f_1f_2 ,
\]
the peer's blueprint's unified form. It is printed here, where both branches
are in scope, and not inside either endpoint lemma, each of which fixes its own
$\epsilon$ and displays its own specialization.
\end{remark}
